\documentclass[12pt,a4paper]{report}
\usepackage[utf8]{inputenc}
\usepackage[turkish]{babel}
\usepackage{amsmath, amssymb, amsthm}
\usepackage{geometry}
\geometry{a4paper, margin=2cm}
\usepackage{tcolorbox}
\usepackage{hyperref}

\title{\textbf{Seriler ve Diziler\\ (Séries et Suites)}}
\author{Ders Notları ve Çözümler}
\date{Université Galatasaray}

\begin{document}
\shorthandoff{=}
\maketitle
\tableofcontents

\chapter{Chapitre 1: Diziler (Suites Numériques)}

\section{TD 0 - Rappels sur les Suites}

\subsection*{Exemple 1.4: [Sınırlı Diziler (Suites Bornées)]}

\begin{tcolorbox}[colback=blue!5!white,colframe=blue!75!black,title=Hatırlatma (Rappel): Sınırlı Diziler]
Bir $(u_n)$ dizisi için:
\begin{itemize}
    \item \textbf{Üstten Sınırlı (Majorée):} Her $n$ için $u_n \le M$ olacak şekilde bir $M \in \mathbb{R}$ varsa.
    \item \textbf{Alttan Sınırlı (Minorée):} Her $n$ için $u_n \ge m$ olacak şekilde bir $m \in \mathbb{R}$ varsa.
    \item \textbf{Sınırlı (Bornée):} Hem alttan hem üstten sınırlıysa (yani $|u_n| \le K$).
\end{itemize}
Trigonometrik fonksiyonlar her zaman sınırlıdır: $-1 \le \sin(x) \le 1$ ve $-1 \le \cos(x) \le 1$.
\end{tcolorbox}

\textbf{Soru:}
Montrer que la suite $u_n = 2 \sin n - 3 \cos(n^2)$ est bornée.

\begin{proof}[Résolution]
Par l'inégalité triangulaire, on a :
\[ |u_n| = |2 \sin n - 3 \cos(n^2)| \le |2 \sin n| + |-3 \cos(n^2)| \]
\[ |u_n| \le 2 |\sin n| + 3 |\cos(n^2)| \]
Or, on sait que pour tout $x \in \mathbb{R}$, $|\sin x| \le 1$ et $|\cos x| \le 1$. En appliquant ceci, on obtient :
\[ |u_n| \le 2(1) + 3(1) = 5 \]
Ainsi, pour tout $n \in \mathbb{N}$, on a $|u_n| \le 5$. Ceci montre que la suite est \textbf{bornée}.
\end{proof}

\subsection*{Exemple 1.6 et 1.7: [Monotonluk (Suites Monotones)]}

\begin{tcolorbox}[colback=green!5!white,colframe=green!60!black,title=Hatırlatma (Rappel): Monotonluk Testleri]
Bir dizinin artan (croissante) veya azalan (décroissante) olduğunu kanıtlamak için iki ana yöntem vardır:
\begin{itemize}
    \item \textbf{Fark Yöntemi:} $u_{n+1} - u_n$ farkına bakılır. $>0$ ise artan, $<0$ ise azalandır. (En genel yöntem).
    \item \textbf{Oran Yöntemi:} Tüm terimler pozitifse ($u_n > 0$), $\frac{u_{n+1}}{u_n}$ oranına bakılır. $>1$ ise artan, $<1$ ise azalandır. (Özellikle faktöriyel ve üslü ifadelerde çok işe yarar).
\end{itemize}
\end{tcolorbox}

\textbf{Soru 1 (Exemple 1.6):} Étudier la monotonie de $u_n = \frac{n}{n+1}$ pour $n \ge 0$.\\
\textbf{Soru 2 (Exemple 1.7):} Étudier la monotonie de $u_n = \frac{2^n}{n!}$ pour $n \ge 1$.

\begin{proof}[Résolution]
\textbf{Partie 1 (Méthode de la différence) :}
Calculons $u_{n+1} - u_n$ :
\[ u_{n+1} - u_n = \frac{n+1}{n+2} - \frac{n}{n+1} \]
En réduisant au même dénominateur :
\[ = \frac{(n+1)^2 - n(n+2)}{(n+1)(n+2)} = \frac{n^2+2n+1 - (n^2+2n)}{(n+1)(n+2)} = \frac{1}{(n+1)(n+2)} \]
Pour tout $n \ge 0$, le dénominateur est strictement positif, donc $u_{n+1} - u_n > 0$.
La suite est donc \textbf{strictement croissante}.

\textbf{Partie 2 (Méthode du quotient) :}
Puisque tous les termes sont strictement positifs ($u_n > 0$), on peut calculer le quotient $\frac{u_{n+1}}{u_n}$ :
\[ \frac{u_{n+1}}{u_n} = \frac{\frac{2^{n+1}}{(n+1)!}}{\frac{2^n}{n!}} = \frac{2^{n+1}}{(n+1)!} \times \frac{n!}{2^n} \]
En simplifiant les puissances et les factorielles :
\[ = \frac{2^n \cdot 2}{n! \cdot (n+1)} \times \frac{n!}{2^n} = \frac{2}{n+1} \]
Pour $n \ge 1$, on a $n+1 \ge 2$, d'où :
\[ \frac{2}{n+1} \le \frac{2}{2} = 1 \]
Ainsi, $\frac{u_{n+1}}{u_n} \le 1$. La suite est donc \textbf{décroissante}.
\end{proof}

\subsection*{Exemple 1.9: [Limiti Doğrudan Kanıtlama ($\epsilon - N$ Tanımı)]}

\begin{tcolorbox}[colback=purple!5!white,colframe=purple!75!black,title=Hatırlatma (Rappel): Limitin Formel Tanımı]
Dizinin limitinin $L$ olması ($u_n \to L$), şu formel anlama gelir:
Öyle bir $N$ (eşik değeri) bulunmalıdır ki, bu $N$'den sonraki her $n$ adımı için dizi elemanı ile $L$ arasındaki mesafe istediğimiz kadar küçük bir $\epsilon > 0$ sayısından daha az olsun.
Matematiksel olarak: $\forall \epsilon > 0, \exists N \in \mathbb{N} \implies (n \ge N \implies |u_n - L| < \epsilon)$.
\end{tcolorbox}

\textbf{Soru:}
Montrons directement que $\lim_{n \to +\infty} \frac{1}{n} = 0$.

\begin{proof}[Démonstration]
Soit $\epsilon > 0$. Nous cherchons un entier $N \in \mathbb{N}$ tel que pour tout $n \ge N$, on ait :
\[ \left| \frac{1}{n} - 0 \right| < \epsilon \]
Ceci équivaut à :
\[ \frac{1}{n} < \epsilon \iff n > \frac{1}{\epsilon} \]
Il suffit de choisir un entier $N > \frac{1}{\epsilon}$ (par exemple $N = \lfloor \frac{1}{\epsilon} \rfloor + 1$).
Alors, pour tout $n \ge N$, on a bien :
\[ \frac{1}{n} \le \frac{1}{N} < \epsilon \]
Ce qui prouve que $\lim_{n \to +\infty} \frac{1}{n} = 0$.
\end{proof}

\subsection*{Exemple 1.18 et 1.20: [Sıkıştırma Teoremi (Théorème des Deux Gendarmes)]}

\begin{tcolorbox}[colback=orange!5!white,colframe=orange!75!black,title=Hatırlatma (Rappel): Sıkıştırma Teoremi (Sandviç Kuralı)]
Dizilerde limit hesaplarken bazen dizinin doğrudan limitini almak çok zordur (örneğin faktöriyel veya trigonometrik karmaşalar içeriyorsa).
Eğer $u_n$ dizisini, aynı limite sahip iki dizi ($v_n$ ve $w_n$) arasına sıkıştırabilirsek ($v_n \le u_n \le w_n$), ortadaki dizi de mecburen o limite gider.
(Fransızlar buna İki Jandarma (Gendarme) Teoremi der: Sağdan ve soldan iki jandarma ortadaki suçluyu karakola ($L$ limitine) götürür).
\end{tcolorbox}

\textbf{Soru 1 (Exemple 1.18):} Calculer $\lim_{n \to +\infty} \frac{3 + \sin(n^2)}{n}$.\\
\textbf{Soru 2 (Exemple 1.20):} Montrer par le théorème des deux gendarmes que $\lim_{n \to +\infty} \frac{2^n}{n!} = 0$.

\begin{proof}[Résolution]
\textbf{Partie 1 :}
Soit $u_n = \frac{3 + \sin(n^2)}{n}$.
On sait que pour tout $n \in \mathbb{N}$, $-1 \le \sin(n^2) \le 1$. En ajoutant 3 :
\[ 2 \le 3 + \sin(n^2) \le 4 \]
En divisant par $n > 0$ :
\[ \frac{2}{n} \le \frac{3 + \sin(n^2)}{n} \le \frac{4}{n} \]
Or, $\lim_{n \to +\infty} \frac{2}{n} = 0$ et $\lim_{n \to +\infty} \frac{4}{n} = 0$.
D'après le théorème des deux gendarmes, on conclut que :
\[ \lim_{n \to +\infty} u_n = 0 \]

\textbf{Partie 2 :}
Soit $u_n = \frac{2^n}{n!}$.
Pour $n \ge 4$, on peut écrire :
\[ n! = 1 \cdot 2 \cdot 3 \cdot 4 \cdots n \]
Cette factorielle contient $(n-3)$ termes qui sont supérieurs ou égaux à 4 :
\[ n! \ge (1 \cdot 2 \cdot 3) \cdot (4 \cdot 4 \cdots 4) = 6 \cdot 4^{n-3} \]
On peut donc majorer la suite :
\[ 0 \le \frac{2^n}{n!} \le \frac{2^n}{6 \cdot 4^{n-3}} \]
En simplifiant le terme de droite :
\[ \frac{2^n}{6 \cdot 4^{n-3}} = \frac{2^n \cdot 4^3}{6 \cdot 4^n} = \frac{64}{6} \left(\frac{2}{4}\right)^n = \frac{32}{3} \left(\frac{1}{2}\right)^n \]
On obtient l'encadrement suivant :
\[ 0 \le u_n \le \frac{32}{3} \left(\frac{1}{2}\right)^n \]
Comme $1/2 \in ]-1, 1[$, la suite géométrique $(1/2)^n$ tend vers 0. Par le théorème des deux gendarmes, on a :
\[ \lim_{n \to +\infty} \frac{2^n}{n!} = 0 \]
\end{proof}

\subsection*{Exemple 1.23: [Monoton ve Sınırlı Dizilerin Yakınsaklığı]}

\begin{tcolorbox}[colback=cyan!5!white,colframe=cyan!75!black,title=Hatırlatma (Rappel): Monotonluk ve Sınırlılık (Convergence des suites monotones)]
Çok önemli bir teorem:
\begin{itemize}
    \item Bir dizi \textbf{artan (croissante)} ve \textbf{üstten sınırlıysa (majorée)}, kesinlikle \textbf{yakınsaktır}.
    \item Bir dizi \textbf{azalan (décroissante)} ve \textbf{alttan sınırlıysa (minorée)}, kesinlikle \textbf{yakınsaktır}.
\end{itemize}
Dikkat: Sadece sınırlı olmak yetmez (örn: $(-1)^n$), sadece monoton olmak da yetmez (örn: $n$). İkisi birleşince yakınsaklığı garanti eder!
\end{tcolorbox}

\textbf{Soru:}
Considérons la suite $u_n = 2 - \frac{1}{n+1}$. Montrer qu'elle est croissante, majorée et par conséquent convergente.

\begin{proof}[Résolution]
Montrons d'abord que la suite est croissante en étudiant $u_{n+1} - u_n$ :
\[ u_{n+1} - u_n = \left(2 - \frac{1}{n+2}\right) - \left(2 - \frac{1}{n+1}\right) \]
\[ = \frac{1}{n+1} - \frac{1}{n+2} = \frac{1}{(n+1)(n+2)} > 0 \]
La suite est donc \textbf{croissante}.

Montrons ensuite qu'elle est majorée :
Pour tout $n \ge 0$, $\frac{1}{n+1} > 0$. Par conséquent :
\[ 2 - \frac{1}{n+1} < 2 \]
La suite est donc \textbf{majorée} par 2.
D'après le théorème de la limite monotone (toute suite croissante et majorée est convergente), la suite $(u_n)$ est \textbf{convergente}. Sa limite est :
\[ \lim_{n \to +\infty} \left(2 - \frac{1}{n+1}\right) = 2 - 0 = 2 \]
\end{proof}

\subsection*{Exemple 1.32: [Cauchy Dizileri (Suites de Cauchy)]}

\begin{tcolorbox}[colback=brown!5!white,colframe=brown!75!black,title=Hatırlatma (Rappel): Cauchy Dizisi]
Bir dizinin Cauchy dizisi olması, terimler ilerledikçe kendi aralarındaki mesafenin (ardışık olmalarına gerek yok, rastgele iki terim $u_n, u_m$) sıfıra yaklaşması demektir.
Reel sayılarda (Dans $\mathbb{R}$), bir dizinin yakınsak olması ile Cauchy dizisi olması birbirine denktir (Critère de Cauchy). 
Formel tanım: $\forall \epsilon > 0, \exists N \in \mathbb{N} \implies (m, n \ge N \implies |u_n - u_m| < \epsilon)$.
\end{tcolorbox}

\textbf{Soru:}
Considérons la suite $u_n = \frac{1}{n}$. Montrons qu'elle est de Cauchy.

\begin{proof}[Démonstration]
Soit $\epsilon > 0$. Prenons deux entiers $n, m \ge N$.
Évaluons la différence $|u_n - u_m|$ :
\[ |u_n - u_m| = \left| \frac{1}{n} - \frac{1}{m} \right| \]
Par l'inégalité triangulaire :
\[ \left| \frac{1}{n} - \frac{1}{m} \right| \le \frac{1}{n} + \frac{1}{m} \]
Comme $n \ge N$ et $m \ge N$, on a :
\[ \frac{1}{n} + \frac{1}{m} \le \frac{1}{N} + \frac{1}{N} = \frac{2}{N} \]
Nous voulons que cette quantité soit strictement inférieure à $\epsilon$ :
\[ \frac{2}{N} < \epsilon \implies N > \frac{2}{\epsilon} \]
Il suffit de choisir un entier $N > \frac{2}{\epsilon}$. Ainsi, pour tout $n, m \ge N$, on a $|u_n - u_m| < \epsilon$.
La suite $(1/n)$ est donc une \textbf{suite de Cauchy}.
\end{proof}

\chapter{Chapitre 2: Seriler (Séries Numériques)}

\section{Serilerin Tanımı ve Kısmi Toplamlar (Sommes Partielles)}

\begin{tcolorbox}[colback=green!5!white,colframe=green!60!black,title=Hatırlatma: Dizi ve Seri Arasındaki Fark]
\begin{itemize}
    \item \textbf{Dizi (Suite) $u_n$:} Sayıların virgülle ardışık dizilmesidir. Örn: $1, \frac{1}{2}, \frac{1}{4}, \dots$
    \item \textbf{Seri (Série) $\sum u_n$:} Bu dizinin elemanlarının sonsuza kadar \textbf{toplanmasıdır}. Örn: $1 + \frac{1}{2} + \frac{1}{4} + \dots$
\end{itemize}
Bir serinin yakınsak olması demek, bu sonsuz toplamın sonlu bir sayıya eşit olması demektir!
\end{tcolorbox}

\section{Temel Seri Testleri (Critères de Convergence)}

\begin{tcolorbox}[colback=red!5!white,colframe=red!75!black,title=Çok Önemli: Genel Terim Testi (Condition Nécessaire)]
Bir serinin $\sum u_n$ yakınsak olabilmesi için \textbf{ZORUNLU} şart şudur:
\[ \lim_{n \to \infty} u_n = 0 \]
Eğer dizinin kendisi sıfıra gitmiyorsa, o sayıları sonsuza kadar toplarsak sonuç asla sonlu çıkmaz! 
\textit{Uyarı: Bu sadece bir ön şarttır. Limit 0 çıksa bile seri ıraksak olabilir (Örn: Harmonik Seri $\sum \frac{1}{n}$).}
\end{tcolorbox}

\section{TD 1 - Premières Séries Numériques}

\subsection*{Exercice 1: [Toplam Sembolü ve İndis Değiştirme (Sommes et Changement d'Indice)]}

\begin{tcolorbox}[colback=purple!5!white,colframe=purple!75!black,title=Hatırlatma (Rappel): Toplam Sembolü Özellikleri]
Toplam sembolü ($\sum$) ile işlem yaparken bazı temel kuralları hatırlayalım:
\begin{itemize}
    \item \textbf{Doğrusallık (Linéarité):} $\sum (a \cdot u_k + b \cdot v_k) = a \sum u_k + b \sum v_k$. Yani sabit katsayıları dışarı alabilir, toplamları ayırabiliriz.
    \item \textbf{Gauss Formülü:} İlk $n$ tam sayının toplamı: $\sum_{k=1}^n k = \frac{n(n+1)}{2}$.
    \item \textbf{Sabit Terim Toplamı:} $c$ bir sabit olmak üzere $\sum_{k=1}^n c = c \times n$.
    \item \textbf{İndis Değiştirme (Changement d'indice):} Toplamın başladığı noktayı değiştirmek için değişkeni öteliyoruz. Yeni bir değişken atamak (örn. $j = k + a$) alt sınır, üst sınır ve içerideki fonksiyonu aynı anda bu yeni değişkene göre güncellemeyi gerektirir.
\end{itemize}
\end{tcolorbox}

\textbf{Soru:}
(a) Calculer $\sum_{k=1}^n (2k - 1)$.\\
(b) Effectuer un changement d'indice dans $\sum_{k=0}^n \frac{1}{k+2}$ afin de l'écrire sous la forme d'une somme dont l'indice commence à 2.

\begin{proof}[Résolution]
\textbf{Partie (a) :}
Par linéarité de la somme, on a :
\[ \sum_{k=1}^n (2k - 1) = 2 \sum_{k=1}^n k - \sum_{k=1}^n 1 \]
En utilisant la formule de Gauss ($\sum_{k=1}^n k = \frac{n(n+1)}{2}$) et le fait que $\sum_{k=1}^n 1 = n$ :
\[ = 2 \left( \frac{n(n+1)}{2} \right) - n = n(n+1) - n = n^2 \]
Ainsi, on obtient :
\[ \sum_{k=1}^n (2k - 1) = n^2 \]

\textbf{Partie (b) :}
On a la somme initiale : 
\[ \sum_{k=0}^n \frac{1}{k+2} \]
Effectuons le changement d'indice $j = k + 2$. Ainsi, $k = j - 2$.
Les nouvelles bornes deviennent :
- Pour $k = 0 \implies j = 2$
- Pour $k = n \implies j = n + 2$

Le terme général devient : $\frac{1}{(j-2)+2} = \frac{1}{j}$.
En remplaçant le tout, on a :
\[ \sum_{k=0}^n \frac{1}{k+2} = \sum_{j=2}^{n+2} \frac{1}{j} \]
L'indice de sommation étant une variable muette, on peut le renommer en $k$ :
\[ \sum_{k=2}^{n+2} \frac{1}{k} \]
\end{proof}

\subsection*{Exercice 2: [Geometrik Seriler (Séries Géométriques)]}

\begin{tcolorbox}[colback=teal!5!white,colframe=teal!75!black,title=Hatırlatma (Rappel): Geometrik Seri ve Kısmi Toplam]
\textbf{Geometrik Seri:} Her terimi bir önceki terimin sabit bir $q$ (raison) sayısıyla çarpılmasıyla elde edilen serilerdir: $\sum q^n = 1 + q + q^2 + q^3 + \dots$ \\
\textbf{Kısmi Toplam (Somme Partielle) Formülü:} İlk $N+1$ terimin toplamı ($S_N$):
\[ S_N = \sum_{n=0}^N q^n = \frac{1 - q^{N+1}}{1 - q} \quad (q \neq 1 \text{ için}) \]
Bu serinin sonsuza kadar yakınsak (convergente) olabilmesi için $N \to \infty$ iken $q^{N+1}$ teriminin sıfıra gitmesi gerekir, ki bu da yalnızca $|q| < 1$ olmasıyla mümkündür.
\end{tcolorbox}

\textbf{Soru:}
Soit $q \in \mathbb{R} \setminus \{1\}$. Montrer que la série géométrique $\sum_{n=0}^{+\infty} q^n$ converge si et seulement si $|q| < 1$. Dans ce cas, calculer sa somme.

\begin{proof}[Résolution]
Considérons la somme partielle $S_N$ :
\[ S_N = \sum_{n=0}^N q^n = 1 + q + q^2 + \dots + q^N \]
En multipliant par $q$ :
\[ q S_N = q + q^2 + q^3 + \dots + q^{N+1} \]
En soustrayant la deuxième équation de la première :
\[ S_N - q S_N = 1 - q^{N+1} \implies S_N = \frac{1 - q^{N+1}}{1 - q} \]
Pour étudier la convergence, regardons la limite quand $N \to +\infty$ :
\[ \lim_{N \to +\infty} S_N = \lim_{N \to +\infty} \frac{1 - q^{N+1}}{1 - q} \]
Le comportement dépend de $q^{N+1}$ :
\begin{itemize}
    \item \textbf{Si $|q| < 1$ :} $\lim_{N \to +\infty} q^{N+1} = 0$, donc $S_N \to \frac{1}{1 - q}$. La série \textbf{converge}.
    \item \textbf{Si $q > 1$ :} $\lim_{N \to +\infty} q^{N+1} = +\infty$, donc la limite diverge vers $+\infty$. La série \textbf{diverge}.
    \item \textbf{Si $q \le -1$ :} $q^{N+1}$ n'admet pas de limite, donc la série \textbf{diverge}.
\end{itemize}
La série converge \textbf{si et seulement si} $|q| < 1$, et sa somme est :
\[ \sum_{n=0}^{+\infty} q^n = \frac{1}{1-q} \]
\end{proof}

\subsection*{Exercice 3: [Genel Terim Testi (Condition Nécessaire de Convergence)]}

\begin{tcolorbox}[colback=red!5!white,colframe=red!75!black,title=Hatırlatma (Rappel): Yakınsaklık İçin Gerekli Şart]
Bir $\sum u_n$ serisinin yakınsak olabilmesi için genel terimi olan $u_n$'nin mutlaka 0'a gitmesi gerekir: $\lim_{n \to \infty} u_n = 0$. 
Eğer bu limit 0'dan farklıysa veya limit yoksa, hiç işlem yapmadan seriye \textbf{"ıraksaktır" (divergente)} deriz. Buna Fransızca'da \textbf{divergence grossière} (kaba ıraksaklık) denir. 
\textit{Ancak dikkat:} Limit 0'a gidiyorsa bu tek başına serinin yakınsak olduğunu kanıtlamaz! (Bkz. sorunun devamı).
\end{tcolorbox}

\textbf{Soru:}
Pour chacune des séries suivantes, étudier la limite du terme général et préciser ce que l'on peut en déduire concernant la convergence :
\[ \sum_{n=1}^{+\infty} \frac{n+2}{n+1}, \quad \sum_{n=1}^{+\infty} \cos\left(\frac{1}{n}\right), \quad \sum_{n=1}^{+\infty} \frac{1}{n^2} \]
La condition $u_n \to 0$ est-elle suffisante pour assurer la convergence de la série $\sum u_n$ ?

\begin{proof}[Résolution]
Étudions la limite du terme général $u_n$ dans chaque cas :

\textbf{1ère série :} $u_n = \frac{n+2}{n+1}$
\[ \lim_{n \to +\infty} \frac{n(1 + \frac{2}{n})}{n(1 + \frac{1}{n})} = \frac{1+0}{1+0} = 1 \]
Comme la limite n'est pas 0 ($1 \neq 0$), la condition nécessaire n'est pas vérifiée. 
\textbf{Conclusion :} La série est \textbf{grossièrement divergente}.

\textbf{2ème série :} $u_n = \cos\left(\frac{1}{n}\right)$
Puisque $1/n \to 0$ et que le cosinus est continu en 0 :
\[ \lim_{n \to +\infty} \cos\left(\frac{1}{n}\right) = \cos(0) = 1 \]
Comme $1 \neq 0$, la condition nécessaire n'est pas vérifiée.
\textbf{Conclusion :} La série est \textbf{grossièrement divergente}.

\textbf{3ème série :} $u_n = \frac{1}{n^2}$
\[ \lim_{n \to +\infty} \frac{1}{n^2} = 0 \]
\textbf{Conclusion :} La condition nécessaire est satisfaite, la série peut converger ou diverger (il se trouve qu'elle converge car c'est une série de Riemann avec $\alpha=2>1$).

\textbf{Réponse à la question :}
Non, la condition $u_n \to 0$ n'est \textbf{pas suffisante} pour assurer la convergence de la série. Un contre-exemple classique est la \textbf{série harmonique} $\sum_{n=1}^{+\infty} \frac{1}{n}$ où $u_n \to 0$ mais la série diverge.
\end{proof}

\subsection*{Exercice 4: [Geometrik Seri Uygulamaları (Calculs de Séries Géométriques)]}

\textbf{Soru:}
Déterminer la nature des séries suivantes et calculer leur somme lorsqu'elles convergent :
\[ \sum_{n=0}^{+\infty} \left(\frac{3}{4}\right)^n, \quad \sum_{n=0}^{+\infty} \left(-\frac{1}{3}\right)^n, \quad \sum_{n=0}^{+\infty} \left(-\frac{5}{4}\right)^n \]

\begin{proof}[Résolution]
On sait que la série géométrique $\sum q^n$ converge si et seulement si $|q| < 1$, de somme $\frac{\text{premier terme}}{1-q}$.

\textbf{1.} $u_n = \left(\frac{3}{4}\right)^n$. On a $q = \frac{3}{4}$. Puisque $|q| < 1$, la série \textbf{converge}.
La somme est :
\[ \sum_{n=0}^{+\infty} \left(\frac{3}{4}\right)^n = \frac{1}{1 - \frac{3}{4}} = \frac{1}{1/4} = 4 \]

\textbf{2.} $u_n = \left(-\frac{1}{3}\right)^n$. On a $q = -\frac{1}{3}$. Puisque $|q| = \frac{1}{3} < 1$, la série \textbf{converge}.
La somme est :
\[ \sum_{n=0}^{+\infty} \left(-\frac{1}{3}\right)^n = \frac{1}{1 - (-\frac{1}{3})} = \frac{1}{4/3} = \frac{3}{4} \]

\textbf{3.} $u_n = \left(-\frac{5}{4}\right)^n$. On a $q = -\frac{5}{4}$. Puisque $|q| = \frac{5}{4} \ge 1$, la condition n'est pas satisfaite.
La série \textbf{diverge}.
\end{proof}

\subsection*{Exercice 5: [Ondalık Yazılımlar ve Geometrik Seriler (Écritures décimales)]}

\begin{tcolorbox}[colback=orange!5!white,colframe=orange!75!black,title=Hatırlatma (Rappel): Devirli Ondalık Sayılar]
Lisede gördüğümüz $0.\overline{3}$ veya $0.\overline{9}$ gibi devirli ondalık (décimal périodique) sayıların ardında yatan asıl mantık serilerdir. Bir sayının ondalık açılımı aslında 10'un negatif kuvvetlerinin bir toplamıdır.
Örneğin: $0.xyz = \frac{x}{10^1} + \frac{y}{10^2} + \frac{z}{10^3}$
Devirli bir kısmı sonsuza kadar gidiyorsa, bu sonsuz toplam (seri) bir geometrik seri oluşturur.
\end{tcolorbox}

\textbf{Soru:}
(1) Écrire $0.3333\dots$ sous la forme d'une série géométrique. Montrer que $0.33\dots = \frac{1}{3}$.\\
(2) En utilisant la même méthode, déterminer la valeur exacte de $0.999999\dots$

\begin{proof}[Résolution]
\textbf{Partie (1) :}
On peut écrire $0.3333\dots$ comme une somme de puissances de 10 :
\[ 0.3333\dots = \frac{3}{10} + \frac{3}{100} + \frac{3}{1000} + \dots = \sum_{n=1}^{+\infty} \frac{3}{10^n} \]
Ceci est une série géométrique de raison $q = \frac{1}{10}$. On met 3 en facteur :
\[ = 3 \sum_{n=1}^{+\infty} \left(\frac{1}{10}\right)^n \]
Puisque $|q| = 1/10 < 1$, la série converge. Son premier terme (pour $n=1$) est $1/10$.
\[ \sum_{n=1}^{+\infty} \left(\frac{1}{10}\right)^n = \frac{1/10}{1 - 1/10} = \frac{1/10}{9/10} = \frac{1}{9} \]
Donc, $0.3333\dots = 3 \times \frac{1}{9} = \frac{1}{3}$.

\textbf{Partie (2) :}
De même pour $0.9999\dots$ :
\[ 0.9999\dots = \frac{9}{10} + \frac{9}{100} + \dots = \sum_{n=1}^{+\infty} \frac{9}{10^n} = 9 \sum_{n=1}^{+\infty} \left(\frac{1}{10}\right)^n \]
D'après le calcul précédent, la somme géométrique vaut $1/9$.
\[ 0.9999\dots = 9 \times \frac{1}{9} = 1 \]
La valeur exacte est bien $1$.
\end{proof}

\subsection*{Exercice 6: [Teleskopik Seriler (Séries Télescopiques)]}

\begin{tcolorbox}[colback=cyan!5!white,colframe=cyan!75!black,title=Hatırlatma (Rappel): Teleskopik Seri]
Teleskopik seriler, genel terimi $u_n = v_{n} - v_{n+1}$ (veya tam tersi) şeklinde yazılabilen serilerdir. Böyle bir serinin kısmi toplamını ($S_N$) açtığımızda, aradaki tüm terimler birbirini sadeleştirir (Domino etkisi gibi).
Sadece baştaki ve sondaki terimler hayatta kalır:
$S_N = (v_1 - v_2) + (v_2 - v_3) + \dots + (v_N - v_{N+1}) = v_1 - v_{N+1}$.
Sonrasında $N \to \infty$ iken limite bakılarak serinin değeri bulunur.
\end{tcolorbox}

\textbf{Soru:}
On considère la série $\sum_{n=2}^{+\infty} \left( \frac{1}{\sqrt{n-1}} - \frac{2}{\sqrt{n}} + \frac{1}{\sqrt{n+1}} \right)$.
Calculer la somme partielle $S_N$. En déduire que la série converge et calculer sa somme.

\begin{proof}[Résolution]
Étudions le terme général $u_n = \frac{1}{\sqrt{n-1}} - \frac{2}{\sqrt{n}} + \frac{1}{\sqrt{n+1}}$.
On peut séparer le terme central en deux :
\[ u_n = \left( \frac{1}{\sqrt{n-1}} - \frac{1}{\sqrt{n}} \right) - \left( \frac{1}{\sqrt{n}} - \frac{1}{\sqrt{n+1}} \right) \]
Posons $v_n = \frac{1}{\sqrt{n-1}} - \frac{1}{\sqrt{n}}$. On a alors $v_{n+1} = \frac{1}{\sqrt{n}} - \frac{1}{\sqrt{n+1}}$.
Ainsi, $u_n = v_n - v_{n+1}$. La série est donc \textbf{télescopique}.

Calculons la somme partielle $S_N$ pour $n \ge 2$ :
\[ S_N = \sum_{n=2}^N (v_n - v_{n+1}) = (v_2 - v_3) + (v_3 - v_4) + \dots + (v_N - v_{N+1}) \]
Tous les termes intermédiaires s'annulent, il reste :
\[ S_N = v_2 - v_{N+1} \]
Avec $v_2 = 1 - \frac{1}{\sqrt{2}}$ et $v_{N+1} = \frac{1}{\sqrt{N}} - \frac{1}{\sqrt{N+1}}$, on obtient :
\[ S_N = 1 - \frac{1}{\sqrt{2}} - \frac{1}{\sqrt{N}} + \frac{1}{\sqrt{N+1}} \]
En passant à la limite quand $N \to +\infty$, on a $\frac{1}{\sqrt{N}} \to 0$ et $\frac{1}{\sqrt{N+1}} \to 0$.
Donc :
\[ \lim_{N \to +\infty} S_N = 1 - \frac{1}{\sqrt{2}} \]
La série \textbf{converge} et sa somme vaut $1 - \frac{1}{\sqrt{2}}$.
\end{proof}

\subsection*{Exercice 7: [Geometrik Serilerin Kombinasyonu (Combinaison)]}

\textbf{Soru:}
On considère la série $\sum_{n=0}^{+\infty} \left( \frac{3}{2^n} + 5 \left(-\frac{2}{7}\right)^n \right)$.
Montrer qu'elle converge. Calculer sa somme.

\begin{proof}[Résolution]
Par linéarité, la somme de deux séries convergentes est convergente. Étudions les deux parties :
\[ \sum_{n=0}^{+\infty} \frac{3}{2^n} \quad \text{et} \quad \sum_{n=0}^{+\infty} 5 \left(-\frac{2}{7}\right)^n \]

\textbf{Première partie :} $3 \sum_{n=0}^{+\infty} \left(\frac{1}{2}\right)^n$
C'est une série géométrique avec $q = 1/2$. Comme $|q| < 1$, elle converge.
Sa somme est : $3 \times \frac{1}{1 - 1/2} = 3 \times 2 = 6$.

\textbf{Deuxième partie :} $5 \sum_{n=0}^{+\infty} \left(-\frac{2}{7}\right)^n$
C'est une série géométrique avec $q = -2/7$. Comme $|q| < 1$, elle converge.
Sa somme est : $5 \times \frac{1}{1 - (-2/7)} = 5 \times \frac{1}{9/7} = 5 \times \frac{7}{9} = \frac{35}{9}$.

La série globale est la somme de deux séries convergentes, elle est donc \textbf{convergente}.
Sa somme totale est :
\[ S = 6 + \frac{35}{9} = \frac{54}{9} + \frac{35}{9} = \frac{89}{9} \]
\end{proof}

\subsection*{Exercice 8: [İntegral Testi ve Bertrand Serileri (Critère Intégral)]}

\begin{tcolorbox}[colback=brown!5!white,colframe=brown!75!black,title=Hatırlatma (Rappel): İntegral Testi ve Limit Karşılaştırması]
\textbf{Limit Karşılaştırması:} $u_n$ ve $v_n$ pozitif iki seri olsun. Eğer $\lim_{n \to \infty} \frac{u_n}{v_n} = 0$ ise, bu $u_n$'nin $v_n$'den çok daha hızlı 0'a gittiği (yani $u_n \ll v_n$) anlamına gelir. \newline
\textbf{İntegral Testi (Critère intégral de Cauchy):} Sürekli, pozitif ve \textbf{azalan (décroissante)} bir $f(x)$ fonksiyonu için, $u_n = f(n)$ serisinin karakteri (yakınsak/ıraksak), $\int_a^\infty f(x) dx$ genelleştirilmiş integralinin karakteriyle \textbf{birebir aynıdır}. İntegral yakınsıyorsa seri de yakınsar, ıraksıyorsa seri de ıraksar.
\end{tcolorbox}

\textbf{Soru:}
On considère les séries $\sum_{n=2}^{+\infty} \frac{1}{n \ln n}$ et $\sum_{n=2}^{+\infty} \frac{1}{n(\ln n)^2}$. Montrer que leurs termes généraux tendent vers 0 plus vite que $\frac{1}{n}$, puis étudier la nature de chacune de ces séries.

\begin{proof}[Résolution]
\textbf{Partie 1 : Vitesse de décroissance}
Comparons les termes généraux avec $1/n$ en calculant le quotient :
Pour la première série :
\[ \lim_{n \to +\infty} \frac{\frac{1}{n \ln n}}{\frac{1}{n}} = \lim_{n \to +\infty} \frac{1}{\ln n} = 0 \]
Pour la deuxième série :
\[ \lim_{n \to +\infty} \frac{\frac{1}{n (\ln n)^2}}{\frac{1}{n}} = \lim_{n \to +\infty} \frac{1}{(\ln n)^2} = 0 \]
Dans les deux cas, le quotient tend vers 0. Ainsi, les termes généraux tendent vers 0 \textbf{plus vite} que $1/n$ ($u_n = o(1/n)$).

\textbf{Partie 2 : Nature des séries par le critère intégral}
Les fonctions $f(x) = \frac{1}{x \ln x}$ et $g(x) = \frac{1}{x (\ln x)^2}$ sont continues, positives et \textbf{décroissantes} sur $[2, +\infty[$.

\textbf{1ère série :} Intégrale $\int_2^X \frac{1}{x \ln x} dx$
On pose $u = \ln x$, donc $du = \frac{dx}{x}$.
\[ \int_{\ln 2}^{\ln X} \frac{1}{u} du = [\ln |u|]_{\ln 2}^{\ln X} = \ln(\ln X) - \ln(\ln 2) \]
Quand $X \to +\infty$, $\ln(\ln X) \to +\infty$. L'intégrale diverge, donc la série \textbf{diverge}.

\textbf{2ème série :} Intégrale $\int_2^X \frac{1}{x (\ln x)^2} dx$
Avec le même changement de variable $u = \ln x$ :
\[ \int_{\ln 2}^{\ln X} \frac{1}{u^2} du = \left[ -\frac{1}{u} \right]_{\ln 2}^{\ln X} = \frac{1}{\ln 2} - \frac{1}{\ln X} \]
Quand $X \to +\infty$, $1/\ln X \to 0$. L'intégrale converge vers $1/\ln 2$. Donc la série \textbf{converge}.
\end{proof}

\subsection*{Exercice 9: [Geometri Uygulaması - Flocon de Koch (Koch Kar Tanesi)]}

\textbf{Soru:}
(Une application : le flocon de Koch) On part d'un triangle équilatéral d'aire 1. À la première étape, on ajoute 3 triangles équilatéraux, chacun d'aire $\frac{1}{9}$. À chaque nouvelle étape, le nombre de triangles ajoutés est multiplié par 4 et l'aire de chacun de ces triangles est divisée par 9. On note $A_n$ l'aire totale obtenue après $n$ étapes.
(1) Calculer $A_1$, $A_2$ et $A_3$.\\
(2) Montrer que $A_n = 1 + \frac{1}{3} \sum_{k=0}^{n-1} \left(\frac{4}{9}\right)^k$.\\
(3) Calculer $\lim_{n \to +\infty} A_n$.

\begin{proof}[Résolution]
\textbf{Partie (1) :}
L'aire initiale est $A_0 = 1$.
À chaque étape $n$, le nombre de nouveaux triangles est multiplié par 4 (soit $3 \times 4^{n-1}$ triangles), et l'aire de chacun est divisée par 9 (soit $1/9^n$).

- \textbf{Pour } $A_1$ : On ajoute 3 triangles d'aire $1/9$. 
Aire ajoutée = $3 \times \frac{1}{9} = \frac{1}{3}$. Donc $A_1 = 1 + \frac{1}{3} = \frac{4}{3}$.
- \textbf{Pour } $A_2$ : On ajoute $3 \times 4 = 12$ triangles d'aire $\frac{1}{81}$.
Aire ajoutée = $12 \times \frac{1}{81} = \frac{4}{27}$. Donc $A_2 = \frac{4}{3} + \frac{4}{27} = \frac{40}{27}$.
- \textbf{Pour } $A_3$ : On ajoute $12 \times 4 = 48$ triangles d'aire $\frac{1}{729}$.
Aire ajoutée = $48 \times \frac{1}{729} = \frac{16}{243}$. Donc $A_3 = \frac{40}{27} + \frac{16}{243} = \frac{376}{243}$.

\textbf{Partie (2) :}
À l'étape $k$, l'aire totale ajoutée est $E_k = (3 \times 4^{k-1}) \times \frac{1}{9^k} = \frac{1}{3} \left(\frac{4}{9}\right)^{k-1}$.
L'aire totale $A_n$ est la somme de l'aire initiale et des aires ajoutées jusqu'à l'étape $n$ :
\[ A_n = 1 + \sum_{k=1}^n \frac{1}{3} \left(\frac{4}{9}\right)^{k-1} \]
En effectuant le changement d'indice $j = k - 1$, la somme varie de $j=0$ à $n-1$ :
\[ A_n = 1 + \frac{1}{3} \sum_{j=0}^{n-1} \left(\frac{4}{9}\right)^j \]
En renommant $j$ en $k$, on obtient :
\[ A_n = 1 + \frac{1}{3} \sum_{k=0}^{n-1} \left(\frac{4}{9}\right)^k \]

\textbf{Partie (3) :}
Quand $n \to +\infty$, la somme $\sum_{k=0}^{n-1} (4/9)^k$ est la somme partielle d'une série géométrique de raison $q = 4/9 < 1$, qui converge vers $\frac{1}{1 - 4/9} = \frac{9}{5}$.
Ainsi, $\lim_{n \to +\infty} A_n = 1 + \frac{1}{3} \times \frac{9}{5} = 1 + \frac{3}{5} = \frac{8}{5}$.
L'aire totale du flocon de Koch est $\frac{8}{5}$.
\end{proof}

\section{TD 3 - Critères de Convergence (Karşılaştırma ve Yakınsaklık Testleri)}

\subsection*{Exercice 1: [Karşılaştırma Testi (Critère de Comparaison)]}

\begin{tcolorbox}[colback=blue!5!white,colframe=blue!75!black,title=Hatırlatma (Rappel): Eşlenik ile Çarpma ve Karşılaştırma]
\textbf{Eşlenik (Quantité Conjuguée):} Köklü ifadelerin farkı durumunda, ifadeyi kökten kurtarmak için eşleniği ile çarparız: $(\sqrt{A} - \sqrt{B}) = \frac{A - B}{\sqrt{A} + \sqrt{B}}$.
\textbf{Karşılaştırma (Comparaison):} Pozitif terimli serilerde, $u_n \le v_n$ ise ve büyük olan $\sum v_n$ yakınsıyorsa, küçük olan $\sum u_n$ de mecburen yakınsar.
\end{tcolorbox}

\textbf{Soru:}
Étudier la nature de la série $\sum_{n \ge 1} \left( \sqrt{n^3+5} - \sqrt{n^3+2} \right)$. Rationaliser le terme général et le comparer à une série de Riemann.

\begin{proof}[Résolution]
Le terme général est $u_n = \sqrt{n^3+5} - \sqrt{n^3+2}$. Multiplions par la quantité conjuguée :
\[ u_n = \frac{(\sqrt{n^3+5} - \sqrt{n^3+2})(\sqrt{n^3+5} + \sqrt{n^3+2})}{\sqrt{n^3+5} + \sqrt{n^3+2}} = \frac{3}{\sqrt{n^3+5} + \sqrt{n^3+2}} \]
Cherchons un équivalent quand $n \to +\infty$. Les termes dominants sont $n^3$.
\[ \sqrt{n^3+5} \sim n^{3/2} \quad \text{et} \quad \sqrt{n^3+2} \sim n^{3/2} \]
Le dénominateur est donc équivalent à $2n^{3/2}$.
\[ u_n \sim \frac{3}{2n^{3/2}} \]
La série de référence est $\sum \frac{1}{n^{3/2}}$, qui est une série de Riemann convergente car $\alpha = 3/2 > 1$.
Par le critère de comparaison par équivalents pour les séries à termes positifs, la série $\sum u_n$ \textbf{converge}.
\end{proof}

\subsection*{Exercice 2: [Asimptotik Denklik (Comparaison par Équivalents)]}

\begin{tcolorbox}[colback=purple!5!white,colframe=purple!75!black,title=Hatırlatma (Rappel): Polinomlarda Denklik]
$n \to \infty$ durumunda, rasyonel fonksiyonlarda (polinom bölü polinom) her zaman \textbf{en yüksek dereceli} terimler baskındır. Diğer küçük dereceli terimleri ve sabitleri limite giderken görmezden gelebiliriz. 
İki pozitif seri için $u_n \sim v_n$ ise, bu iki seri "kardeş" olur, yani biri yakınsıyorsa diğeri de yakınsar, biri ıraksıyorsa diğeri de ıraksar.
\end{tcolorbox}

\textbf{Soru:}
Étudier la nature des séries suivantes. Dans chaque cas, déterminer un équivalent simple du terme général et comparer avec une série de Riemann :
(1) $\sum_{n \ge 1} \frac{3n^2+2n+1}{n^4+1}$ \qquad (2) $\sum_{n \ge 1} \frac{1}{\sqrt{n^2+3n}}$

\begin{proof}[Résolution]
\textbf{Partie (1) :} $u_n = \frac{3n^2+2n+1}{n^4+1}$.
En ne gardant que les termes de plus haut degré :
\[ u_n \sim \frac{3n^2}{n^4} = \frac{3}{n^2} \]
La série $\sum \frac{1}{n^2}$ est une série de Riemann avec $\alpha = 2 > 1$, donc convergente.
Par le critère d'équivalence, la série \textbf{converge}.

\textbf{Partie (2) :} $u_n = \frac{1}{\sqrt{n^2+3n}}$.
De même :
\[ u_n \sim \frac{1}{\sqrt{n^2}} = \frac{1}{n} \]
La série harmonique $\sum \frac{1}{n}$ est divergente (Riemann avec $\alpha = 1 \le 1$).
Par le critère d'équivalence, la série \textbf{diverge}.
\end{proof}

\subsection*{Exercice 3: [Riemann Serileri (Séries de Riemann)]}

\textbf{Soru:}
(1) Déterminer la nature de $\sum_{n \ge 1} \frac{1}{n^{3/4}}$.\\
(2) Étudier la série $\sum_{n \ge 1} \frac{1}{n^{(n+1)/n}}$ en remarquant que $\frac{1}{n^{(n+1)/n}} = \frac{1}{n} \frac{1}{n^{1/n}}$.

\begin{proof}[Résolution]
\textbf{Partie (1) :} $\sum \frac{1}{n^{3/4}}$.
Il s'agit d'une série de Riemann avec $\alpha = 3/4 \le 1$. La série \textbf{diverge}.

\textbf{Partie (2) :} $u_n = \frac{1}{n^{(n+1)/n}}$.
Réécrivons le terme général :
\[ u_n = \frac{1}{n^{1 + 1/n}} = \frac{1}{n} \times \frac{1}{n^{1/n}} \]
Cherchons la limite de $n^{1/n}$ en passant à l'exponentielle :
\[ n^{1/n} = \exp\left( \frac{\ln n}{n} \right) \]
Puisque $\lim_{n \to +\infty} \frac{\ln n}{n} = 0$, on a $\lim_{n \to +\infty} n^{1/n} = e^0 = 1$.
Ainsi, $u_n \sim \frac{1}{n}$.
Puisque la série harmonique diverge, par le critère d'équivalence, la série \textbf{diverge}.
\end{proof}

\subsection*{Exercice 4: [İntegral Testi (Critère Intégral)]}

\textbf{Soru:}
Pour $\alpha \in \mathbb{R}$, on considère la série $\sum_{n \ge 2} \frac{1}{n(\ln n)^\alpha}$.
(1) On pose $f(x) = \frac{1}{x(\ln x)^\alpha}$. Montrer que, pour $\alpha > 0$, la fonction $f$ est positive et décroissante sur $[2, +\infty[$.
(2) À l'aide du changement de variable $t = \ln x$, étudier l'intégrale $\int_2^{+\infty} \frac{dx}{x(\ln x)^\alpha}$.
(3) En déduire que la série converge si et seulement si $\alpha > 1$.

\begin{proof}[Résolution]
\textbf{Partie (1) :} Pour $x \ge 2$, on a $x > 0$ et $\ln x > 0$. Puisque $\alpha > 0$, la fonction $f(x) = \frac{1}{x(\ln x)^\alpha}$ est bien positive.
Calculons la dérivée $f'(x)$ :
\[ f'(x) = -1 \cdot (x(\ln x)^\alpha)^{-2} \cdot \left[ 1 \cdot (\ln x)^\alpha + x \cdot \alpha (\ln x)^{\alpha-1} \frac{1}{x} \right] \]
\[ = - \frac{(\ln x)^\alpha + \alpha (\ln x)^{\alpha-1}}{(x(\ln x)^\alpha)^2} \]
Pour $x \ge 2$ et $\alpha > 0$, le numérateur est positif. Avec le signe moins, $f'(x) < 0$. $f$ est donc \textbf{décroissante}.

\textbf{Partie (2) :} Intégrale $I = \int_2^{X} \frac{dx}{x(\ln x)^\alpha}$.
Posons $t = \ln x$, d'où $dt = \frac{dx}{x}$.
\[ I = \int_{\ln 2}^{\ln X} \frac{dt}{t^\alpha} \]
- Si $\alpha \neq 1$, $I = \left[ \frac{t^{1-\alpha}}{1-\alpha} \right]_{\ln 2}^{\ln X}$. Quand $X \to +\infty$, cette intégrale converge si $1-\alpha < 0 \iff \alpha > 1$.
- Si $\alpha = 1$, $I = [\ln|t|]_{\ln 2}^{\ln X} = \ln(\ln X) - \ln(\ln 2)$, qui diverge.
L'intégrale converge si et seulement si $\alpha > 1$.

\textbf{Partie (3) :} D'après le critère intégral (pour une fonction positive décroissante), la série et l'intégrale ont la même nature.
La série \textbf{converge si et seulement si} $\alpha > 1$.
\end{proof}

\subsection*{Exercice 5: [D'Alembert (Oran) Testi (Critère de d'Alembert)]}

\begin{tcolorbox}[colback=teal!5!white,colframe=teal!75!black,title=Hatırlatma (Rappel): D'Alembert Kriteri]
D'Alembert (veya Ratio/Oran) Testi, özellikle faktöriyel (!) ve üslü terimler ($a^n$) içeren serilerde kullanılır.
Pozitif bir seri için $L = \lim_{n \to \infty} \frac{u_{n+1}}{u_n}$ hesaplanır:
\begin{itemize}
    \item $L < 1 \implies$ Seri \textbf{yakınsak (convergente)}.
    \item $L > 1 \implies$ Seri \textbf{ıraksak (divergente)}.
    \item $L = 1 \implies$ \textbf{Kararsız durum} (Test sonuç vermez, Riemann vb. başka bir test gerekir).
\end{itemize}
Hatırlatma: $\lim_{n \to \infty} \left(1 + \frac{1}{n}\right)^n = e \approx 2.718$.
\end{tcolorbox}

\textbf{Soru:}
Étudier la nature des séries suivantes avec d'Alembert:
(1) $\sum_{n \ge 1} \frac{n!}{n^n}$ \qquad (2) $\sum_{n \ge 1} \frac{2^n n!}{n^n}$

\begin{proof}[Résolution]
\textbf{Partie (1) :} $u_n = \frac{n!}{n^n}$. Calculons le rapport de d'Alembert :
\[ \frac{u_{n+1}}{u_n} = \frac{(n+1)!}{(n+1)^{n+1}} \times \frac{n^n}{n!} = \frac{(n+1)}{(n+1)(n+1)^n} n^n = \left(\frac{n}{n+1}\right)^n = \frac{1}{(1 + \frac{1}{n})^n} \]
Quand $n \to +\infty$, $(1+1/n)^n \to e$. Donc la limite est $L = \frac{1}{e}$.
Comme $e \approx 2.718$, $L < 1$. D'après la règle de d'Alembert, la série \textbf{converge}.

\textbf{Partie (2) :} $u_n = \frac{2^n n!}{n^n}$.
Le rapport donne :
\[ \frac{u_{n+1}}{u_n} = \frac{2^{n+1}}{2^n} \times \frac{(n+1)!}{(n+1)^{n+1}} \times \frac{n^n}{n!} = 2 \times \frac{1}{(1 + \frac{1}{n})^n} \]
La limite est $L = \frac{2}{e}$.
Comme $2 < e$, on a $L < 1$. La série \textbf{converge}.
\end{proof}

\subsection*{Exercice 6: [Kriter Seçimi (Choix d'un critère)]}

\textbf{Soru:}
Étudier la nature des séries suivantes en choisissant un critère adapté. Justifier le choix.
(1) $\sum_{n \ge 2} \frac{\ln n}{n^{3/2}}$ (On pourra utiliser ln n $\le$ n$^{1/4}$)\\
(2) $\sum_{n \ge 2} \frac{1}{n^{1+1/\ln n}}$

\begin{proof}[Résolution]
\textbf{Partie (1) :} $u_n = \frac{\ln n}{n^{3/2}}$.
Utilisons la majoration $\ln n \le n^{1/4}$ pour $n$ assez grand :
\[ u_n \le \frac{n^{1/4}}{n^{3/2}} = \frac{1}{n^{3/2 - 1/4}} = \frac{1}{n^{5/4}} \]
La série de Riemann $\sum \frac{1}{n^{5/4}}$ converge ($\alpha = 5/4 > 1$).
Par le critère de comparaison, la série $\sum u_n$ \textbf{converge}.

\textbf{Partie (2) :} $u_n = \frac{1}{n^{1+1/\ln n}} = \frac{1}{n \cdot n^{1/\ln n}}$.
Simplifions le terme $n^{1/\ln n}$ :
\[ n^{1/\ln n} = \exp\left( \frac{1}{\ln n} \ln n \right) = e^1 = e \]
Ainsi, $u_n = \frac{1}{e \cdot n}$.
La série harmonique $\sum \frac{1}{n}$ diverge. Comme $u_n$ est proportionnel à la série harmonique, la série \textbf{diverge}.
\end{proof}

\subsection*{Exercice 7: [Büyüme Hızları (Vitesses de croissance)]}

\textbf{Soru:}
Soient $a > 1$ et $p > 0$.
(1) À l'aide du critère de d'Alembert, étudier la convergence des séries $\sum \frac{n^p}{a^n}$ et $\sum \frac{a^n}{n!}$.
(2) En déduire que $\lim \frac{n^p}{a^n} = 0$ et $\lim \frac{a^n}{n!} = 0$. Interpréter.

\begin{proof}[Résolution]
\textbf{Partie (1) :}
- Pour $u_n = \frac{n^p}{a^n}$ :
\[ \frac{u_{n+1}}{u_n} = \frac{(n+1)^p}{a^{n+1}} \times \frac{a^n}{n^p} = \left(\frac{n+1}{n}\right)^p \frac{1}{a} = \left(1 + \frac{1}{n}\right)^p \frac{1}{a} \xrightarrow{n \to \infty} \frac{1}{a} \]
Comme $a > 1$, on a $1/a < 1$, donc la série \textbf{converge}.
- Pour $v_n = \frac{a^n}{n!}$ :
\[ \frac{v_{n+1}}{v_n} = \frac{a^{n+1}}{(n+1)!} \times \frac{n!}{a^n} = \frac{a}{n+1} \xrightarrow{n \to \infty} 0 \]
Comme $0 < 1$, la série \textbf{converge}.

\textbf{Partie (2) :}
D'après la condition nécessaire de convergence, le terme général d'une série convergente doit tendre vers 0.
Par conséquent : $\lim_{n \to +\infty} \frac{n^p}{a^n} = 0$ et $\lim_{n \to +\infty} \frac{a^n}{n!} = 0$.
\textbf{Interprétation :} L'exponentielle croît plus vite que toute puissance ($n^p \ll a^n$), et la factorielle croît plus vite que l'exponentielle ($a^n \ll n!$).
\end{proof}

\section{TD 4 - Séries Numériques à Termes Quelconques}

\subsection*{Exercice 1: [Mutlak Yakınsaklık (Convergence absolue)]}

\begin{tcolorbox}[colback=blue!5!white,colframe=blue!75!black,title=Hatırlatma (Rappel): Mutlak Yakınsaklık]
Bir seri negatif terimler de içerebilir. Eğer serinin mutlak değerini aldığımızda elde ettiğimiz pozitif terimli seri yakınsıyorsa ($\sum |u_n|$ convergente), orijinal seriye \textbf{mutlak yakınsak (absolument convergente)} denir.
Çok önemli teorem: \textbf{Mutlak yakınsak her seri, normalde de yakınsaktır.}
\[ \sum |u_n| \text{ converge } \implies \sum u_n \text{ converge} \]
\end{tcolorbox}

\textbf{Soru:}
On considère les séries suivantes :
\[ \sum_{n \ge 1} \frac{(-1)^n}{n^2}, \quad \sum_{n \ge 1} \frac{\sin n}{n^2}, \quad \sum_{n \ge 1} \frac{(-1)^n(2 + \cos n)}{n^2} \]
(1) Lesquelles de ces séries sont alternées?\\
(2) Étudier la convergence absolue de chacune de ces séries.\\
(3) En déduire leur convergence.

\begin{proof}[Résolution]
\textbf{Partie (1) :} Une série est dite alternée si son terme général change de signe à chaque itération (de la forme $(-1)^n a_n$ avec $a_n \ge 0$).
- La première série est \textbf{alternée} car $\frac{1}{n^2} \ge 0$.
- La deuxième série n'est \textbf{pas alternée} car le signe de $\sin n$ change de manière irrégulière.
- La troisième série est \textbf{alternée} car pour tout $n$, $\cos n \ge -1$, donc $2 + \cos n \ge 1 > 0$.

\textbf{Partie (2) :} Étudions la convergence absolue en prenant la valeur absolue du terme général :
- Pour la première : $\left| \frac{(-1)^n}{n^2} \right| = \frac{1}{n^2}$. C'est une série de Riemann avec $\alpha = 2 > 1$, donc elle converge. La série est \textbf{absolument convergente}.
- Pour la deuxième : $\left| \frac{\sin n}{n^2} \right| = \frac{|\sin n|}{n^2} \le \frac{1}{n^2}$. Par le critère de comparaison, comme $\sum \frac{1}{n^2}$ converge, la série est \textbf{absolument convergente}.
- Pour la troisième : $\left| \frac{(-1)^n(2 + \cos n)}{n^2} \right| = \frac{2 + \cos n}{n^2} \le \frac{3}{n^2}$. De même, par comparaison avec une série de Riemann convergente, elle est \textbf{absolument convergente}.

\textbf{Partie (3) :} D'après le théorème fondamental, toute série absolument convergente est convergente. On en déduit que \textbf{les trois séries convergent}.
\end{proof}

\subsection*{Exercice 2: [Leibniz Testi (Critère de Leibniz)]}

\begin{tcolorbox}[colback=purple!5!white,colframe=purple!75!black,title=Hatırlatma (Rappel): Alterne Seriler İçin Leibniz Kriteri]
Alterne bir seri ($(-1)^n a_n$, burada $a_n \ge 0$) şu iki şartı sağlıyorsa kesinlikle \textbf{yakınsaktır}:
1. Dizi $a_n$ azalan (décroissante) olmalıdır.
2. Limit sıfıra gitmelidir ($\lim_{n \to \infty} a_n = 0$).
Eğer seri Leibniz ile yakınsıyor ancak mutlak değerce ıraksıyorsa, buna \textbf{koşullu yakınsak (semi-convergente)} denir.
\end{tcolorbox}

\textbf{Soru:}
Étudier la nature des séries suivantes et pour chacune (1) vérifier les hypothèses du critère de Leibniz, (2) en déduire la convergence, et déterminer si la convergence est absolue.
\[ \sum_{n \ge 1} \frac{(-1)^{n+1} n}{n^2 + 1}, \quad \sum_{n \ge 1} \frac{(-1)^{n+1}}{\sqrt{n}} \]

\begin{proof}[Résolution]
\textbf{Première série :} $u_n = (-1)^{n+1} a_n$ avec $a_n = \frac{n}{n^2+1}$.
\begin{itemize}
    \item \textbf{Vérification de Leibniz :}
    La limite est claire : $\lim_{n \to +\infty} \frac{n}{n^2+1} = \lim_{n \to +\infty} \frac{1}{n} = 0$.
    Pour la monotonie, posons $f(x) = \frac{x}{x^2+1}$. Sa dérivée est $f'(x) = \frac{1(x^2+1) - x(2x)}{(x^2+1)^2} = \frac{1-x^2}{(x^2+1)^2}$. Pour $x \ge 1$, $f'(x) \le 0$, donc la suite $a_n$ est \textbf{décroissante}.
    Les hypothèses de Leibniz sont vérifiées, donc la série \textbf{converge}.
    \item \textbf{Convergence absolue :}
    $|u_n| = \frac{n}{n^2+1} \sim \frac{1}{n}$. La série harmonique $\sum \frac{1}{n}$ diverge. Par le critère d'équivalence, la série des valeurs absolues diverge. 
    \textbf{Conclusion :} La série est \textbf{semi-convergente} (convergente mais pas absolument).
\end{itemize}

\textbf{Deuxième série :} $u_n = (-1)^{n+1} a_n$ avec $a_n = \frac{1}{\sqrt{n}}$.
\begin{itemize}
    \item \textbf{Vérification de Leibniz :}
    $\lim_{n \to +\infty} \frac{1}{\sqrt{n}} = 0$. La suite $a_n$ est trivialement \textbf{décroissante} car la fonction racine est croissante au dénominateur. La série \textbf{converge}.
    \item \textbf{Convergence absolue :}
    $|u_n| = \frac{1}{\sqrt{n}} = \frac{1}{n^{1/2}}$. C'est une série de Riemann avec $\alpha = 1/2 \le 1$, donc elle diverge.
    \textbf{Conclusion :} La série est \textbf{semi-convergente}.
\end{itemize}
\end{proof}

\subsection*{Exercice 3: [Alterne Seriyi Tanıma (Reconnaître une série alternée)]}

\textbf{Soru:}
On considère la série $\sum_{n \ge 0} \frac{\sin(n\pi/2)}{n+1}$.
(1) Calculer $\sin(n\pi/2)$ pour $n = 0, 1, 2, 3, 4, 5, 6, 7$.\\
(2) Montrer que les termes d'indice pair sont nuls.\\
(3) Montrer que, pour $n = 2k + 1$, on a $\sin((2k + 1)\pi/2) = (-1)^k$.\\
(4) En déduire que la série peut s'écrire $\sum_{k \ge 0} \frac{(-1)^k}{2k + 2}$.\\
(5) Montrer que cette série converge.\\
(6) La convergence est-elle absolue?

\begin{proof}[Résolution]
\textbf{Partie (1) :} Valeurs successives de $\sin(n\pi/2)$ pour $n = 0, \dots, 7$ :
$0, 1, 0, -1, 0, 1, 0, -1$.

\textbf{Partie (2) :} Pour un indice pair $n = 2k$, on a :
$\sin\left(\frac{2k\pi}{2}\right) = \sin(k\pi) = 0$ pour tout $k \in \mathbb{N}$. Les termes d'indice pair sont bien nuls.

\textbf{Partie (3) :} Pour un indice impair $n = 2k + 1$, on a :
$\sin\left(\frac{(2k+1)\pi}{2}\right) = \sin\left(k\pi + \frac{\pi}{2}\right)$.
D'après les formules trigonométriques, $\sin(k\pi + \pi/2) = \cos(k\pi) = (-1)^k$.

\textbf{Partie (4) :} Puisque les termes pairs sont nuls, la somme infinie se réduit aux termes impairs. Remplaçons $n$ par $2k+1$ :
Le dénominateur devient $n+1 = (2k+1)+1 = 2k+2$.
La série s'écrit donc : $\sum_{k=0}^{+\infty} \frac{(-1)^k}{2k+2}$.

\textbf{Partie (5) :} On a une série alternée $\sum (-1)^k a_k$ avec $a_k = \frac{1}{2k+2}$.
La suite $(a_k)$ est manifestement décroissante et tend vers 0. D'après le \textbf{critère de Leibniz}, la série converge.

\textbf{Partie (6) :} Étudions la convergence absolue :
$|u_k| = \frac{1}{2k+2} \sim \frac{1}{2k}$.
La série $\sum \frac{1}{2k} = \frac{1}{2} \sum \frac{1}{k}$ est divergente (série harmonique). La convergence \textbf{n'est pas absolue} (elle est semi-convergente).
\end{proof}

\subsection*{Exercice 4: [Cauchy Kök Testi (Critère de Cauchy)]}

\begin{tcolorbox}[colback=teal!5!white,colframe=teal!75!black,title=Hatırlatma (Rappel): Cauchy Kök Testi (Critère de Cauchy)]
Eğer bir serinin genel teriminin tamamı $n$. kuvvet formundaysa (yani $u_n = (a_n)^n$ gibi), Cauchy kök testi uygulanır:
\[ L = \lim_{n \to \infty} \sqrt[n]{|u_n|} \]
- $L < 1 \implies$ Seri mutlak yakınsak.
- $L > 1 \implies$ Seri ıraksak.
- $L = 1 \implies$ Kararsız (Başka test gerekir).
Dikkat: Limit alınırken ünlü $\lim_{n \to \infty} \left(1 + \frac{x}{n}\right)^n = e^x$ formülü sıkça çıkar.
\end{tcolorbox}

\textbf{Soru:}
Étudier la nature des séries suivantes et pour chacune (1) calculer $\sqrt[n]{|u_n|}$, (2) déterminer sa limite et conclure sur la nature de la série.
\[ \sum_{n \ge 1} \left( \frac{2n+1}{3n+2} \right)^n, \quad \sum_{n \ge 1} \left( \frac{n^2-2}{3n^2} \right)^n, \quad \sum_{n \ge 3} \left( 1 - \frac{2}{n} \right)^{n^2} \]

\begin{proof}[Résolution]
\textbf{Première série :} $u_n = \left( \frac{2n+1}{3n+2} \right)^n$
\begin{itemize}
    \item Racine $n$-ième : $\sqrt[n]{u_n} = \frac{2n+1}{3n+2}$.
    \item Limite : En prenant les termes de plus haut degré, $\lim_{n \to +\infty} \frac{2n}{3n} = \frac{2}{3}$.
    \item Conclusion : Comme $L = 2/3 < 1$, d'après le critère de Cauchy, la série \textbf{converge}.
\end{itemize}

\textbf{Deuxième série :} $u_n = \left( \frac{n^2-2}{3n^2} \right)^n$
\begin{itemize}
    \item Racine $n$-ième : $\sqrt[n]{u_n} = \frac{n^2-2}{3n^2}$.
    \item Limite : $\lim_{n \to +\infty} \frac{n^2}{3n^2} = \frac{1}{3}$.
    \item Conclusion : Comme $L = 1/3 < 1$, d'après le critère de Cauchy, la série \textbf{converge}.
\end{itemize}

\textbf{Troisième série :} $u_n = \left( 1 - \frac{2}{n} \right)^{n^2}$
\begin{itemize}
    \item Racine $n$-ième : $\sqrt[n]{u_n} = \left( \left( 1 - \frac{2}{n} \right)^{n^2} \right)^{1/n} = \left( 1 - \frac{2}{n} \right)^n$.
    \item Limite : D'après la limite classique $\lim_{n \to +\infty} \left(1 + \frac{x}{n}\right)^n = e^x$, on a ici $x = -2$, donc $L = e^{-2} = \frac{1}{e^2}$.
    \item Conclusion : Puisque $e \approx 2.718$, $e^2 > 1$, donc $L = \frac{1}{e^2} < 1$. La série \textbf{converge}.
\end{itemize}
\end{proof}

\subsection*{Exercice 5: [Mutlak Yakınsaklığın Özellikleri (Propriétés de la convergence absolue)]}

\textbf{Soru:}
Soit $\sum_{n \ge 0} u_n$ une série réelle absolument convergente. On pose $u_n^+ = \max(u_n, 0)$ et $u_n^- = \max(-u_n, 0)$.
(1) Montrer que $u_n = u_n^+ - u_n^-$ et $|u_n| = u_n^+ + u_n^-$.\\
(2) Montrer que $0 \le u_n^+ \le |u_n|$ et $0 \le u_n^- \le |u_n|$.\\
(3) En déduire que les séries $\sum u_n^+$ et $\sum u_n^-$ convergent.\\
(4) Montrer que $\sum u_n = \sum u_n^+ - \sum u_n^-$.

\begin{proof}[Résolution]
\textbf{Partie (1) :} Procédons par disjonction des cas.
- Si $u_n \ge 0$ : $\max(u_n, 0) = u_n$ et $\max(-u_n, 0) = 0$. Donc $u_n^+ = u_n$ et $u_n^- = 0$.
On a bien $u_n^+ - u_n^- = u_n - 0 = u_n$ et $u_n^+ + u_n^- = u_n + 0 = u_n = |u_n|$ (car $u_n$ est positif).
- Si $u_n < 0$ : $\max(u_n, 0) = 0$ et $\max(-u_n, 0) = -u_n$. Donc $u_n^+ = 0$ et $u_n^- = -u_n$.
On a bien $u_n^+ - u_n^- = 0 - (-u_n) = u_n$ et $u_n^+ + u_n^- = 0 + (-u_n) = -u_n = |u_n|$ (car $u_n$ est négatif).
Dans les deux cas, les relations sont vérifiées.

\textbf{Partie (2) :} 
Par définition, un maximum avec 0 est toujours positif, donc $u_n^+ \ge 0$ et $u_n^- \ge 0$.
De la relation $|u_n| = u_n^+ + u_n^-$, comme ce sont deux termes positifs, on a trivialement :
$u_n^+ \le |u_n|$ et $u_n^- \le |u_n|$.
Ainsi : $0 \le u_n^+ \le |u_n|$ et $0 \le u_n^- \le |u_n|$.

\textbf{Partie (3) :} 
Par hypothèse, la série $\sum u_n$ est absolument convergente, c'est-à-dire que $\sum |u_n|$ converge.
Puisque $0 \le u_n^+ \le |u_n|$ et $0 \le u_n^- \le |u_n|$, par le \textbf{critère de comparaison} des séries à termes positifs, les séries $\sum u_n^+$ et $\sum u_n^-$ \textbf{convergent}.

\textbf{Partie (4) :} 
Puisque $\sum u_n^+$ et $\sum u_n^-$ sont deux séries convergentes, on peut utiliser la \textbf{linéarité} des séries convergentes :
\[ \sum_{n \ge 0} (u_n^+ - u_n^-) = \sum_{n \ge 0} u_n^+ - \sum_{n \ge 0} u_n^- \]
Or d'après (1), $u_n^+ - u_n^- = u_n$, donc :
\[ \sum_{n \ge 0} u_n = \sum_{n \ge 0} u_n^+ - \sum_{n \ge 0} u_n^- \]
Ce qui achève la démonstration.
\end{proof}

\chapter{Chapitre 3: Fonksiyon Dizileri ve Serileri (Suites et Séries de Fonctions)}

\begin{tcolorbox}[colback=teal!5!white,colframe=teal!75!black,title=Temel Teori: Noktasal ve Düzgün Yakınsaklık (Convergence Simple vs Uniforme)]
Fonksiyon dizileri ($f_n(x)$), kalkülüsün ve analizin en önemli konularından biridir. İki farklı yakınsaklık türü vardır:
\begin{enumerate}
    \item \textbf{Noktasal Yakınsaklık (Convergence Simple - $\xrightarrow{C.S.}$):} Her bir $x \in D$ sabitlenir ve sayısal dizi gibi limitine bakılır: $f(x) = \lim_{n \to \infty} f_n(x)$. Burada $N$ sayısı hem $\epsilon$'a hem de $x$'e bağlıdır ($N(\epsilon, x)$).
    \item \textbf{Düzgün Yakınsaklık (Convergence Uniforme - $\xrightarrow{C.U.}$):} Fonksiyon grafiğinin tamamı $f(x)$ limit fonksiyonunun etrafındaki $\epsilon$ şeridine aynı anda girer. Burada seçilen $N$ sayısı \textbf{bütün $x \in D$ noktaları için ortaktır} ($N(\epsilon)$).
    \[ \|f_n - f\|_{\infty, D} = \sup_{x \in D} |f_n(x) - f(x)| \xrightarrow[n \to +\infty]{} 0 \]
\end{enumerate}
\textbf{Altın Kural (Sürekliliğin Aktarılması):} Eğer her bir $f_n$ fonksiyonu sürekli ise ve $f_n \xrightarrow{C.U.} f$ düzgün yakınsıyorsa, limit fonksiyonu $f$ de \textbf{kesinlikle süreklidir}. Eğer $f$ süreksiz çıkarsa, yakınsaklık \textbf{asla düzgün olamaz}!
\end{tcolorbox}

\section{TD 5 (14 Octobre) - Séries Conditionnelles et Suites de Fonctions}

\subsection*{Exercice 1: [Somme de deux séries (Şartlı Yakınsaklık ve Eşdeğerlik Tehlikesi)]}

\begin{tcolorbox}[colback=red!5!white,colframe=red!75!black,title=Çok Önemli Uyarı: İşareti Değişen Serilerde Eşdeğerlik Kuralı GEÇERSİZDİR!]
Pozitif terimli serilerde $u_n \sim v_n$ ise serilerin karakteri aynıdır. Fakat \textbf{işareti değişen serilerde eşdeğerlik kuralı kullanılamaz!} İki dizi birbirine denk olsa bile biri yakınsak diğeri ıraksak olabilir.
\end{tcolorbox}

\textbf{Soru:}
On pose, pour $n \ge 1$,
\[ u_n = \frac{(-1)^n}{\sqrt{n}}, \quad v_n = \frac{(-1)^n}{\sqrt{n}} + \frac{1}{n} \]
(1) Montrer que la série $\sum u_n$ converge conditionnellement.\\
(2) Étudier la convergence de la série $\sum v_n$.\\
(3) Étudier la convergence de la série $\sum (u_n + v_n)$.\\
(4) Que montre cet exemple sur l'addition de deux séries dont les termes tendent vers zéro ?

\begin{proof}[Résolution]
\textbf{Partie (1) :} Étudions la série $\sum u_n = \sum \frac{(-1)^n}{\sqrt{n}}$.
\begin{itemize}
    \item \textbf{Convergence de la série :} C'est une série alternée de la forme $\sum (-1)^n a_n$ avec $a_n = \frac{1}{\sqrt{n}}$. 
    La suite $(a_n)$ est à termes strictement positifs, strictement décroissante car $n \mapsto \sqrt{n}$ est strictement croissante, et $\lim_{n \to +\infty} a_n = \lim_{n \to +\infty} \frac{1}{\sqrt{n}} = 0$.
    D'après le \textbf{critère spécial des séries alternées (Leibniz)}, la série $\sum u_n$ \textbf{converge}.
    \item \textbf{Convergence absolue :}
    $|u_n| = \left| \frac{(-1)^n}{\sqrt{n}} \right| = \frac{1}{\sqrt{n}} = \frac{1}{n^{1/2}}$.
    Il s'agit d'une série de Riemann d'exposant $\alpha = \frac{1}{2} \le 1$. Par conséquent, la série $\sum |u_n|$ \textbf{diverge}.
\end{itemize}
Puisque la série converge mais ne converge pas absolument, la série $\sum u_n$ est \textbf{conditionnellement convergente} (ou semi-convergente).

\textbf{Partie (2) :} Étudions la série $\sum v_n$.
On remarque que pour tout $n \ge 1$ :
\[ v_n = u_n + \frac{1}{n} \iff \frac{1}{n} = v_n - u_n \]
Supposons par l'absurde que la série $\sum v_n$ converge. 
Puisque $\sum u_n$ converge (d'après la partie 1), la combinaison linéaire de deux séries convergentes serait convergente, ce qui impliquerait que la série :
\[ \sum_{n \ge 1} (v_n - u_n) = \sum_{n \ge 1} \frac{1}{n} \]
convergerait également. Or, $\sum \frac{1}{n}$ est la série harmonique (Riemann avec $\alpha = 1$), qui est notoirement \textbf{divergente} ($\sum_{n=1}^{+\infty} \frac{1}{n} = +\infty$).
Cette contradiction prouve que la série $\sum v_n$ \textbf{diverge}.

\textbf{Partie (3) :} Étudions la somme $\sum (u_n + v_n)$.
On a :
\[ u_n + v_n = \frac{(-1)^n}{\sqrt{n}} + \frac{(-1)^n}{\sqrt{n}} + \frac{1}{n} = \frac{2(-1)^n}{\sqrt{n}} + \frac{1}{n} \]
De même, si $\sum (u_n + v_n)$ convergeait, alors comme la série $\sum \frac{2(-1)^n}{\sqrt{n}} = 2 \sum u_n$ converge par linéarité, la différence :
\[ \sum_{n \ge 1} \left( (u_n + v_n) - 2u_n \right) = \sum_{n \ge 1} \frac{1}{n} \]
convergerait, ce qui est absurde. Donc la série $\sum (u_n + v_n)$ \textbf{diverge}.

\textbf{Partie (4) :} Remarque sur cet exemple :
Cet exemple met en évidence deux propriétés fondamentales :
\begin{enumerate}
    \item On observe que $v_n = \frac{(-1)^n}{\sqrt{n}} \left(1 + \frac{(-1)^n}{\sqrt{n}}\right) \sim \frac{(-1)^n}{\sqrt{n}} = u_n$ quand $n \to +\infty$. Pourtant, $\sum u_n$ converge alors que $\sum v_n$ diverge. Cela prouve rigoureusement que \textbf{la règle des équivalents est en général fausse pour les séries qui ne sont pas à termes de signe constant}.
    \item L'addition d'un terme général qui tend vers zéro ($\frac{1}{n} \to 0$) à une série convergente peut détruire la convergence et rendre la série divergente, car la série de ce terme perturbateur n'est pas sommable.
\end{enumerate}
\end{proof}

\subsection*{Exercice 2: [Convergence absolue et conditionnelle (Combinaisons de Séries)]}

\textbf{Soru:}
Étudier la nature des séries suivantes. Lorsqu'une série converge, préciser si sa convergence est absolue ou conditionnelle.
\[ (a) \sum_{n \ge 1} \frac{(-1)^n}{3n-2}, \quad (b) \sum_{n \ge 1} \left( \frac{(-1)^n}{n} + \frac{1}{n^2} \right), \quad (c) \sum_{n \ge 1} \left( \frac{(-1)^n}{n} + \frac{(-1)^{n+1}}{n} + \frac{1}{n^2} \right) \]
(1) Étudier séparément les séries intervenant dans les expressions (b) et (c).\\
(2) Peut-on obtenir une série absolument convergente en additionnant deux séries conditionnellement convergentes?

\begin{proof}[Résolution]
\textbf{Étude de la série (a) :} $u_n = \frac{(-1)^n}{3n-2}$.
\begin{itemize}
    \item \textbf{Convergence :} C'est une série alternée avec $a_n = \frac{1}{3n-2} > 0$.
    On a $a_{n+1} - a_n = \frac{1}{3n+1} - \frac{1}{3n-2} = \frac{-3}{(3n+1)(3n-2)} < 0$, donc $(a_n)$ est strictement décroissante.
    De plus, $\lim_{n \to +\infty} a_n = 0$.
    Par le critère de Leibniz, la série $\sum u_n$ \textbf{converge}.
    \item \textbf{Convergence absolue :}
    $|u_n| = \frac{1}{3n-2} \sim \frac{1}{3n}$.
    La série $\sum \frac{1}{n}$ est divergente (série harmonique). Par le critère d'équivalence pour les séries à termes strictement positifs, $\sum |u_n|$ \textbf{diverge}.
\end{itemize}
Conclusion pour (a) : La série est \textbf{conditionnellement convergente}.

\textbf{Étude de la série (b) et partie (1) :}
L'expression est la somme de $a_n = \frac{(-1)^n}{n}$ et $b_n = \frac{1}{n^2}$.
\begin{itemize}
    \item $\sum a_n = \sum \frac{(-1)^n}{n}$ converge par le critère de Leibniz (série harmonique alternée), mais ne converge pas absolument car $\sum |a_n| = \sum \frac{1}{n}$ diverge. Elle est donc \textbf{conditionnellement convergente}.
    \item $\sum b_n = \sum \frac{1}{n^2}$ est une série de Riemann avec $\alpha = 2 > 1$, elle est donc \textbf{absolument convergente}.
    \item Par linéarité, la somme de deux séries convergentes converge, donc $\sum \left( \frac{(-1)^n}{n} + \frac{1}{n^2} \right)$ \textbf{converge}.
    \item Étudions la convergence absolue :
    \[ \left| \frac{(-1)^n}{n} + \frac{1}{n^2} \right| = \frac{1}{n} \left| (-1)^n + \frac{1}{n} \right| \sim \frac{1}{n} \]
    car $\left| (-1)^n + \frac{1}{n} \right| \xrightarrow[n \to +\infty]{} 1$.
    Par comparaison par équivalents avec la série harmonique divergente $\sum \frac{1}{n}$, la série des valeurs absolues \textbf{diverge}.
\end{itemize}
Conclusion pour (b) : La série est \textbf{conditionnellement convergente}.

\textbf{Étude de la série (c) et partie (1) :}
Remarquons que $(-1)^{n+1} = -(-1)^n$. Par conséquent :
\[ \frac{(-1)^n}{n} + \frac{(-1)^{n+1}}{n} = \frac{(-1)^n}{n} - \frac{(-1)^n}{n} = 0 \]
Le terme général se réduit donc exactement à :
\[ w_n = \frac{1}{n^2} \]
pour tout $n \ge 1$.
Puisque $\sum \frac{1}{n^2}$ est une série de Riemann convergente avec $\alpha = 2 > 1$, la série (c) \textbf{converge absolument}.

\textbf{Partie (2) :} 
\textbf{Oui}, on peut parfaitement obtenir une série absolument convergente en additionnant deux séries conditionnellement convergentes.
L'exemple (c) en fournit une preuve directe : 
Posons $x_n = \frac{(-1)^n}{n}$ et $y_n = \frac{(-1)^{n+1}}{n} + \frac{1}{n^2}$.
Ces deux séries sont chacune conditionnellement convergentes (par les questions précédentes).
Cependant, leur somme vaut :
\[ x_n + y_n = \frac{1}{n^2} \]
dont la série $\sum (x_n + y_n) = \sum \frac{1}{n^2}$ est \textbf{absolument convergente}.
(De façon encore plus triviale : si $\sum u_n$ est semi-convergente, alors $\sum (-u_n)$ l'est aussi, et leur somme est la série nulle $\sum 0 = 0$, qui est absolument convergente).
\end{proof}

\subsection*{Exercice 3: [Convergence simple et uniforme de $f_n(x) = x^n$ (Rôle du Domaine)]}

\begin{tcolorbox}[colback=blue!5!white,colframe=blue!75!black,title=Hatırlatma: Bir Fonksiyon Dizisinin Düzgün Yakınsaklığı Nasıl İncelenir?]
Bir $(f_n)$ fonksiyon dizisinin $D$ kümesinde $f$'e düzgün yakınsak olup olmadığını incelemek için şu adımlar izlenir:
\begin{enumerate}
    \item Önce her $x \in D$ için noktasal limit $f(x) = \lim_{n \to \infty} f_n(x)$ bulunur.
    \item $f_n$ fonksiyonları sürekli olduğu halde limit $f$ süreksiz çıkarsa $\implies$ \textbf{Düzgün yakınsak DEĞİLDİR}.
    \item Eğer $f$ sürekli ise, hata fonksiyonunun tepe noktasına (supremumuna) bakılır:
    \[ M_n = \sup_{x \in D} |f_n(x) - f(x)| \]
    Eğer $\lim_{n \to \infty} M_n = 0$ ise \textbf{düzgün yakınsaktır}, sıfırdan farklı bir sayı veya sonsuz çıkarsa \textbf{düzgün yakınsak değildir}.
\end{enumerate}
\end{tcolorbox}

\textbf{Soru:}
Pour $n \ge 1$, on considère la fonction $f_n : [0, 1] \to \mathbb{R}$ définie par $f_n(x) = x^n$.
(1) Déterminer la limite simple $f$ de la suite $(f_n)$ sur $[0, 1]$.\\
(2) La convergence est-elle uniforme sur $[0, 1]$ ?\\
(3) Soit $a \in ]0, 1[$. Montrer que $(f_n)$ converge uniformément vers 0 sur $[0, a]$.\\
(4) La suite $(f_n)$ converge-t-elle uniformément vers 0 sur $[0, 1[$ ?\\
(5) Comparer les résultats obtenus sur $[0, 1]$, $[0, 1[$ et $[0, a]$. Que peut-on en conclure sur le rôle du domaine dans la convergence uniforme ?

\begin{proof}[Résolution]
\textbf{Partie (1) : Limite simple sur $[0, 1]$}
Soit $x \in [0, 1]$.
\begin{itemize}
    \item Si $x \in [0, 1[$, on a $|x| < 1$, donc $\lim_{n \to +\infty} x^n = 0$.
    \item Si $x = 1$, pour tout $n \ge 1$, $f_n(1) = 1^n = 1$, donc $\lim_{n \to +\infty} f_n(1) = 1$.
\end{itemize}
La suite $(f_n)$ converge simplement sur $[0, 1]$ vers la fonction $f$ définie par :
\[ f(x) = \begin{cases} 0 & \text{si } x \in [0, 1[ \\ 1 & \text{si } x = 1 \end{cases} \]

\textbf{Partie (2) : Convergence uniforme sur $[0, 1]$}
\begin{itemize}
    \item \textbf{Méthode par la continuité (Théorème fondamental) :}
    Pour tout $n \ge 1$, la fonction $f_n : x \mapsto x^n$ est un polynôme, donc elle est continue sur $[0, 1]$.
    Si la suite $(f_n)$ convergeait uniformément vers $f$ sur $[0, 1]$, alors la fonction limite $f$ serait nécessairement continue sur $[0, 1]$.
    Or, on constate que :
    \[ \lim_{x \to 1^-} f(x) = 0 \neq f(1) = 1 \]
    La fonction $f$ est discontinue au point $x = 1$. Par contraposée du théorème de continuité, la suite $(f_n)$ \textbf{ne converge pas uniformément} vers $f$ sur $[0, 1]$.
    \item \textbf{Vérification directe par la norme uniforme :}
    Pour tout $n \ge 1$ :
    \[ \|f_n - f\|_{\infty, [0, 1]} = \sup_{x \in [0, 1]} |f_n(x) - f(x)| \ge \sup_{x \in [0, 1[} |x^n - 0| = \sup_{x \in [0, 1[} x^n = 1 \]
    Comme cette borne supérieure vaut 1 pour tout $n$, elle ne tend pas vers 0 quand $n \to +\infty$.
\end{itemize}

\textbf{Partie (3) : Convergence uniforme sur $[0, a]$ avec $a \in ]0, 1[$}
Sur l'intervalle $[0, a]$, la limite simple est identiquement nulle : $f(x) = 0$.
Calculons la norme uniforme sur $[0, a]$ :
\[ \|f_n - f\|_{\infty, [0, a]} = \sup_{x \in [0, a]} |x^n - 0| = \sup_{x \in [0, a]} x^n \]
Puisque la fonction $x \mapsto x^n$ est strictement croissante et positive sur $[0, a]$, son maximum est atteint à la borne supérieure $x = a$ :
\[ \sup_{x \in [0, a]} x^n = a^n \]
Comme $a \in ]0, 1[$, on a $\lim_{n \to +\infty} a^n = 0$.
Par conséquent, $\lim_{n \to +\infty} \|f_n - f\|_{\infty, [0, a]} = 0$.
La suite $(f_n)$ \textbf{converge uniformément vers 0 sur $[0, a]$}.

\textbf{Partie (4) : Convergence uniforme sur $[0, 1[$}
Sur $[0, 1[$, la limite simple est $f(x) = 0$.
Calculons la borne supérieure sur $[0, 1[$ :
\[ \|f_n - f\|_{\infty, [0, 1[} = \sup_{x \in [0, 1[} |x^n - 0| = \sup_{x \in [0, 1[} x^n = \lim_{x \to 1^-} x^n = 1 \]
Puisque $\|f_n - f\|_{\infty, [0, 1[} = 1 \not\to 0$ quand $n \to +\infty$, la suite $(f_n)$ \textbf{ne converge pas uniformément vers 0 sur $[0, 1[$}.
\textit{Remarque :} Ici la fonction limite $f=0$ est continue sur $[0, 1[$, ce qui montre bien que la continuité de la limite est une condition nécessaire mais non suffisante pour assurer la convergence uniforme.

\textbf{Partie (5) : Rôle du domaine dans la convergence uniforme}
En comparant les résultats sur $[0, 1]$, $[0, 1[$ et $[0, a]$ :
\begin{itemize}
    \item La convergence est uniforme sur tout segment compact $[0, a] \subset [0, 1[$. On dit qu'il y a \textbf{convergence uniforme sur tout compact} (ou convergence compacte).
    \item En revanche, la convergence n'est uniforme ni sur l'intervalle ouvert $[0, 1[$, ni sur l'intervalle fermé $[0, 1]$.
\end{itemize}
\textbf{Conclusion :} La convergence uniforme dépend fondamentalement de la géométrie de l'ensemble de définition : s'éloigner du bord à problème ($x=1$) permet de récupérer la convergence uniforme, alors que s'approcher arbitrairement de ce bord détruit l'uniformité de la convergence.
\end{proof}

\subsection*{Exercice 4: [Convergence uniforme de $f_n(x) = x^n(1-x)$ (Sınır Etkisinin Yok Edilmesi)]}

\textbf{Soru:}
Pour $n \ge 1$, on considère la fonction $f_n : [0, 1] \to \mathbb{R}$ définie par $f_n(x) = x^n(1-x)$.\\
(1) Déterminer la limite simple de $(f_n)$ sur $[0, 1]$.\\
(2) (En utilisant $1 - x^{n+1} = (1-x)(1 + x + \cdots + x^n)$) Montrer que, pour tout $x \in [0, 1]$,
\[ 0 \le f_n(x) \le \frac{1}{n+1} \]
(3) En déduire que $(f_n)$ converge uniformément vers 0 sur $[0, 1]$.\\
(4) Comparer ce résultat avec celui obtenu pour $g_n(x) = x^n$.

\begin{proof}[Résolution]
\textbf{Partie (1) : Limite simple sur $[0, 1]$}
Soit $x \in [0, 1]$.
\begin{itemize}
    \item Pour $x \in [0, 1[$, $\lim_{n \to +\infty} x^n = 0$, donc $\lim_{n \to +\infty} f_n(x) = 0 \cdot (1-x) = 0$.
    \item Pour $x = 1$, pour tout $n \ge 1$, $f_n(1) = 1^n (1-1) = 0$.
\end{itemize}
La suite $(f_n)$ converge simplement sur $[0, 1]$ vers la fonction identiquement nulle : $f(x) = 0$.

\textbf{Partie (2) : Majoration fondamentale}
Soit $x \in [0, 1]$. Puisque $x \ge 0$ et $1-x \ge 0$, on a immédiatement $f_n(x) = x^n(1-x) \ge 0$.
Utilisons l'identité donnée :
\[ 1 - x^{n+1} = (1-x)(1 + x + x^2 + \cdots + x^n) = \sum_{k=0}^n (1-x)x^k \]
Pour $x \in [0, 1]$, on a pour tout $k \in \{0, 1, \dots, n\}$, $x^n \le x^k$.
Puisque $1-x \ge 0$, on peut multiplier par $(1-x)$ :
\[ x^n (1-x) \le x^k (1-x) \]
En sommant ces inégalités pour $k$ variant de $0$ à $n$ :
\[ \sum_{k=0}^n x^n(1-x) \le \sum_{k=0}^n x^k(1-x) \]
Le membre de gauche contient $(n+1)$ termes identiques égaux à $f_n(x)$ :
\[ (n+1) f_n(x) \le 1 - x^{n+1} \]
Puisque $x \in [0, 1]$, on a $x^{n+1} \ge 0$, d'où $1 - x^{n+1} \le 1$. Il s'ensuit que :
\[ (n+1) f_n(x) \le 1 \implies f_n(x) \le \frac{1}{n+1} \]
On a donc bien établi que pour tout $x \in [0, 1]$ :
\[ 0 \le f_n(x) \le \frac{1}{n+1} \]

\textit{Remarque alternative par dérivée :} On a $f_n'(x) = n x^{n-1}(1-x) - x^n = x^{n-1}[n - (n+1)x]$. Le maximum est atteint en $x_n = \frac{n}{n+1}$, et $f_n(x_n) = \left(\frac{n}{n+1}\right)^n \frac{1}{n+1} \le \frac{1}{n+1}$.

\textbf{Partie (3) : Convergence uniforme}
D'après l'encadrement précédent, pour tout $x \in [0, 1]$ :
\[ |f_n(x) - 0| = f_n(x) \le \frac{1}{n+1} \]
Comme cette majoration est indépendante de $x$, on passe à la borne supérieure sur $[0, 1]$ :
\[ \|f_n - 0\|_{\infty, [0, 1]} = \sup_{x \in [0, 1]} |f_n(x)| \le \frac{1}{n+1} \]
Or $\lim_{n \to +\infty} \frac{1}{n+1} = 0$.
Par le théorème des gendarmes, $\lim_{n \to +\infty} \|f_n - 0\|_{\infty, [0, 1]} = 0$.
La suite $(f_n)$ \textbf{converge uniformément vers 0 sur $[0, 1]$}.

\textbf{Partie (4) : Comparaison avec $g_n(x) = x^n$}
Pour la suite $g_n(x) = x^n$, la convergence n'était pas uniforme sur $[0, 1]$ car les valeurs restaient proches de 1 lorsque $x \to 1^-$ (discontinuité au bord).
Dans le cas de $f_n(x) = x^n(1-x)$, la présence du facteur amortisseur $(1-x)$ force la fonction à s'annuler exactement en $x=1$ : il compense la croissance de $x^n$ et écrase la bosse vers zéro au voisinage de 1. Cela régularise le comportement au bord et transforme une convergence non uniforme en une \textbf{convergence uniforme} sur tout le compact $[0, 1]$.
\end{proof}

\subsection*{Exercice 5: [Suites de Fonctions Classiques (Recherche de Suites Test)]}

\begin{tcolorbox}[colback=purple!5!white,colframe=purple!75!black,title=Sınav Taktiği: Düzgün Yakınsak Olmadığını Göstermek İçin "Suite Test" Bulma]
Eğer bir fonksiyon dizisinin $D$ üzerinde düzgün yakınsak \textbf{olmadığını} ispatlamak istiyorsak, öyle bir $x_n \in D$ dizisi seçeriz ki:
\[ |f_n(x_n) - f(x_n)| \xrightarrow[n \to \infty]{} c \neq 0 \]
Bu dizi genellikle fonksiyonun türevini sıfır yapan maksimum noktasıdır (örneğin $x_n = \frac{1}{n}$, $x_n = n$ veya sınıra yaklaşan bir nokta).
\end{tcolorbox}

\textbf{Soru:}
Pour chacune des suites de fonctions suivantes : (1) déterminer la limite simple sur le domaine indiqué, (2) étudier la convergence uniforme, (3) lorsque la convergence n'est pas uniforme, trouver si possible une suite $(x_n)$ de points du domaine telle que $|f_n(x_n) - f(x_n)|$ ne tende pas vers 0.
\[ (a) f_n(x) = \frac{x}{n}, x \in [0, 1]; \quad (b) f_n(x) = \frac{x}{n}, x \in [0, +\infty[ \]
\[ (c) f_n(x) = \frac{nx}{1 + n^2 x^2}, x \in [0, 1]; \quad (d) f_n(x) = \frac{x^n}{1 + x^n}, x \in [0, 1] \]

\begin{proof}[Résolution]
\textbf{Cas (a) :} $f_n(x) = \frac{x}{n}$ sur $D = [0, 1]$.
\begin{enumerate}
    \item \textbf{Limite simple :} Pour chaque $x \in [0, 1]$ fixé, $\lim_{n \to +\infty} \frac{x}{n} = 0$. Donc la limite simple est $f(x) = 0$.
    \item \textbf{Convergence uniforme :}
    $\|f_n - 0\|_{\infty, [0, 1]} = \sup_{x \in [0, 1]} \frac{x}{n} = \frac{1}{n} \xrightarrow[n \to +\infty]{} 0$.
    La convergence est donc \textbf{uniforme sur $[0, 1]$}.
\end{enumerate}

\textbf{Cas (b) :} $f_n(x) = \frac{x}{n}$ sur $D = [0, +\infty[$.
\begin{enumerate}
    \item \textbf{Limite simple :} Pour tout $x \in [0, +\infty[$ fixé, $\lim_{n \to +\infty} \frac{x}{n} = 0$. Donc $f(x) = 0$.
    \item \textbf{Convergence uniforme :}
    $\|f_n - 0\|_{\infty, [0, +\infty[} = \sup_{x \ge 0} \frac{x}{n} = +\infty \neq 0$.
    La convergence \textbf{n'est pas uniforme sur $[0, +\infty[$}.
    \item \textbf{Suite test :} Posons $x_n = n \in [0, +\infty[$. Alors :
    \[ |f_n(x_n) - f(x_n)| = \frac{n}{n} = 1 \xrightarrow[n \to +\infty]{} 1 \neq 0 \]
\end{enumerate}

\textbf{Cas (c) :} $f_n(x) = \frac{nx}{1 + n^2 x^2}$ sur $D = [0, 1]$.
\begin{enumerate}
    \item \textbf{Limite simple :}
    - Pour $x = 0$, $f_n(0) = 0 \to 0$.\\
    - Pour $x \in ]0, 1]$, $f_n(x) = \frac{nx}{1 + n^2 x^2} \sim \frac{nx}{n^2 x^2} = \frac{1}{nx} \xrightarrow[n \to +\infty]{} 0$.\\
    Donc la limite simple est identiquement nulle sur tout $[0, 1]$ : $f(x) = 0$.
    \item \textbf{Convergence uniforme :}
    Étudions les variations de $f_n$ sur $[0, 1]$ en calculant sa dérivée :
    \[ f_n'(x) = \frac{n(1 + n^2 x^2) - nx(2 n^2 x)}{(1 + n^2 x^2)^2} = \frac{n(1 - n^2 x^2)}{(1 + n^2 x^2)^2} \]
    La dérivée s'annule en $x_n = \frac{1}{n} \in [0, 1]$ pour $n \ge 1$.
    La fonction $f_n$ est strictement croissante sur $[0, 1/n]$ puis strictement décroissante sur $[1/n, 1]$.
    Son maximum absolu est donc atteint en $x_n = \frac{1}{n}$ :
    \[ f_n\left(\frac{1}{n}\right) = \frac{n(1/n)}{1 + n^2 (1/n^2)} = \frac{1}{1 + 1} = \frac{1}{2} \]
    Par conséquent :
    \[ \|f_n - 0\|_{\infty, [0, 1]} = \sup_{x \in [0, 1]} |f_n(x)| = \frac{1}{2} \xrightarrow[n \to +\infty]{} \frac{1}{2} \neq 0 \]
    La convergence \textbf{n'est pas uniforme sur $[0, 1]$}.
    \item \textbf{Suite test :} La suite des extremums $x_n = \frac{1}{n} \in [0, 1]$ donne :
    \[ |f_n(x_n) - f(x_n)| = f_n(1/n) - 0 = \frac{1}{2} \xrightarrow[n \to +\infty]{} \frac{1}{2} \neq 0 \]
\end{enumerate}

\textbf{Cas (d) :} $f_n(x) = \frac{x^n}{1 + x^n}$ sur $D = [0, 1]$.
\begin{enumerate}
    \item \textbf{Limite simple :}
    - Si $x \in [0, 1[$, $x^n \to 0$, d'où $f_n(x) = \frac{x^n}{1+x^n} \to \frac{0}{1+0} = 0$.\\
    - Si $x = 1$, $f_n(1) = \frac{1^n}{1+1^n} = \frac{1}{2}$.\\
    La limite simple est :
    \[ f(x) = \begin{cases} 0 & \text{si } x \in [0, 1[ \\ \frac{1}{2} & \text{si } x = 1 \end{cases} \]
    \item \textbf{Convergence uniforme :}
    Chaque $f_n$ est continue sur $[0, 1]$ (quotient de fonctions continues dont le dénominateur ne s'annule pas car $1+x^n \ge 1$).
    Or, la fonction limite $f$ présente une discontinuité en $x = 1$ car $\lim_{x \to 1^-} f(x) = 0 \neq f(1) = \frac{1}{2}$.
    Par le théorème de conservation de la continuité, la convergence \textbf{ne peut pas être uniforme sur $[0, 1]$}.
    \item \textbf{Suite test :} Choisissons $x_n = \left(\frac{1}{2}\right)^{1/n} \in [0, 1[$.
    Alors $x_n^n = \frac{1}{2}$.
    Puisque $x_n \in [0, 1[$, on a $f(x_n) = 0$.
    On évalue alors :
    \[ |f_n(x_n) - f(x_n)| = \frac{1/2}{1 + 1/2} - 0 = \frac{1/2}{3/2} = \frac{1}{3} \xrightarrow[n \to +\infty]{} \frac{1}{3} \neq 0 \]
\end{enumerate}
\end{proof}

\subsection*{Exercice 6: [Questions de Réflexion (Sous-Ensembles et Segments Compacts)]}

\textbf{Soru:}
Soit $(f_n)$ une suite de fonctions qui converge uniformément vers $f$ sur un ensemble $D$.\\
(1) Si $A \subset D$, la convergence de $(f_n)$ vers $f$ est-elle nécessairement uniforme sur $A$ ? Justifier.\\
(2) Supposons maintenant que $(f_n)$ converge uniformément vers $f$ sur tout segment contenu dans $D$. Peut-on en déduire que $(f_n)$ converge uniformément vers $f$ sur $D$ ?\\
(3) Donner un contre-exemple.

\begin{proof}[Résolution]
\textbf{Partie (1) : Restriction à un sous-ensemble $A \subset D$}
\textbf{Oui}, la convergence est nécessairement uniforme sur tout sous-ensemble $A \subset D$.
\begin{itemize}
    \item \textbf{Démonstration :} Par propriété fondamentale de la borne supérieure, si $A \subset D$, alors :
    \[ \sup_{x \in A} |f_n(x) - f(x)| \le \sup_{x \in D} |f_n(x) - f(x)| \]
    Par hypothèse, $(f_n)$ converge uniformément vers $f$ sur $D$, ce qui signifie que :
    \[ \lim_{n \to +\infty} \sup_{x \in D} |f_n(x) - f(x)| = 0 \]
    Puisque pour tout $n$, $0 \le \sup_{x \in A} |f_n(x) - f(x)|$, le \textbf{théorème des gendarmes} s'applique directement et donne :
    \[ \lim_{n \to +\infty} \sup_{x \in A} |f_n(x) - f(x)| = 0 \]
    Ce qui prouve que la suite converge uniformément sur $A$.
\end{itemize}

\textbf{Partie (2) et (3) : Réciproque et contre-exemple}
\textbf{Non}, on ne peut pas en déduire que $(f_n)$ converge uniformément sur $D$.
La convergence uniforme sur tout segment compact (notée convergence uniforme locale ou convergence compacte) est strictement plus faible que la convergence uniforme globale sur $D$.

\textbf{Contre-exemple (deux illustrations fondamentales) :}
\begin{enumerate}
    \item \textbf{Exemple 1 (Intervalle borné non fermé) :} 
    Considérons $D = [0, 1[$ et $f_n(x) = x^n$.
    - Tout segment contenu dans $D$ est de la forme $[a, b]$ avec $0 \le a \le b < 1$. Sur ce segment, la fonction étant croissante :
    \[ \sup_{x \in [a, b]} |f_n(x) - 0| = b^n \xrightarrow[n \to +\infty]{} 0 \quad (\text{car } b < 1) \]
    La convergence est donc uniforme sur tout segment de $[0, 1[$.
    - Pourtant, sur $D = [0, 1[$, on a :
    \[ \sup_{x \in [0, 1[} |f_n(x) - 0| = \sup_{x \in [0, 1[} x^n = 1 \xrightarrow[n \to +\infty]{} 1 \neq 0 \]
    La convergence n'est pas uniforme sur $D$.
    \item \textbf{Exemple 2 (Domaine non borné) :} 
    Considérons $D = [0, +\infty[$ et $f_n(x) = \frac{x}{n}$.
    - Tout segment de $D$ est inclus dans un intervalle $[0, M]$ avec $M > 0$. Sur ce segment :
    \[ \sup_{x \in [0, M]} \left| \frac{x}{n} \right| = \frac{M}{n} \xrightarrow[n \to +\infty]{} 0 \]
    La convergence est uniforme sur tout segment de $[0, +\infty[$.
    - Pourtant, sur $D = [0, +\infty[$, on a $\sup_{x \ge 0} \frac{x}{n} = +\infty$, donc la convergence n'est pas uniforme sur $D$.
\end{enumerate}
\end{proof}

\section{TD 5 Suite / TD 6 (21 Octobre) - Suites et Séries de Fonctions, Continuité et Cauchy Uniforme}

\begin{tcolorbox}[colback=green!5!white,colframe=green!60!black,title=Temel Teori: Fonksiyon Serilerinde Yakınsaklık Hiyerarşisi (Séries de Fonctions)]
Bir fonksiyon serisi $\sum f_n(x)$ için yakınsaklık türleri arasında şu hiyerarşi geçerlidir:
\[ \text{\textbf{Normal Yakınsaklık}} \implies \text{\textbf{Düzgün Yakınsaklık}} \implies \text{\textbf{Noktasal Yakınsaklık}} \]
\[ \text{\textbf{Normal Yakınsaklık}} \implies \text{\textbf{Mutlak Yakınsaklık}} \implies \text{\textbf{Noktasal Yakınsaklık}} \]
\begin{itemize}
    \item \textbf{Normal Yakınsaklık (Convergence Normale):} $\sum \|f_n\|_{\infty, D} < +\infty$ olmasıdır. En güçlü yakınsaklıktır! Sayısal seriye indirger (\textbf{Weierstrass $M$-testi}).
    \item \textbf{Düzgün Yakınsaklık (Convergence Uniforme):} Kısmi toplamlar dizisinin $S_N(x)$ düzgün yakınsamasıdır. Bu da \textbf{kalan terimin} (reste) sıfıra gitmesine denktir:
    \[ \lim_{N \to +\infty} \sup_{x \in D} |R_N(x)| = 0, \quad \text{où } R_N(x) = \sum_{n=N+1}^{+\infty} f_n(x) \]
    \item \textbf{DİKKAT (Tersleri Yanlıştır!):} Düzgün yakınsak bir seri normal yakınsak olmak zorunda \textbf{değildir} (en meşhur karşıt örnek: $\sum \frac{(-1)^n x}{n}$ serisi).
\end{itemize}
\end{tcolorbox}

\subsection*{Exercice 1: [Convergence simple et uniforme de suites de fonctions]}

\textbf{Soru:}
Pour chacune des suites de fonctions suivantes, déterminer la limite simple, puis étudier la convergence uniforme sur le domaine indiqué.
\[ a) f_n(x) = \frac{x}{1+nx}, x \in [0, 1] \]
\[ b) f_n(x) = \frac{x}{n+x}, x \in [0, +\infty[ \]
\[ c) f_n(x) = \frac{x^2}{n+x^2}, x \in \mathbb{R} \]

\begin{proof}[Résolution]
\textbf{Cas a) :} $f_n(x) = \frac{x}{1+nx}$ sur $D = [0, 1]$.
\begin{itemize}
    \item \textbf{Limite simple :}
    - Pour $x = 0$ : $f_n(0) = 0 \to 0$.\\
    - Pour $x \in ]0, 1]$ : $f_n(x) = \frac{x}{1+nx} \sim \frac{x}{nx} = \frac{1}{n} \xrightarrow[n \to +\infty]{} 0$.\\
    Donc la limite simple est la fonction nulle sur $[0, 1]$ : $f(x) = 0$.
    \item \textbf{Convergence uniforme :}
    Calculons la dérivée de $f_n$ sur $[0, 1]$ :
    \[ f_n'(x) = \frac{1 \cdot (1+nx) - x \cdot n}{(1+nx)^2} = \frac{1}{(1+nx)^2} > 0 \]
    La fonction $f_n$ est strictement croissante sur $[0, 1]$. Son maximum est donc atteint en $x = 1$ :
    \[ \|f_n - 0\|_{\infty, [0, 1]} = \sup_{x \in [0, 1]} f_n(x) = f_n(1) = \frac{1}{1+n} \]
    Puisque $\lim_{n \to +\infty} \frac{1}{1+n} = 0$, on a :
    \[ \lim_{n \to +\infty} \|f_n - f\|_{\infty, [0, 1]} = 0 \]
    La suite $(f_n)$ \textbf{converge uniformément vers 0 sur $[0, 1]$}.
\end{itemize}

\textbf{Cas b) :} $f_n(x) = \frac{x}{n+x}$ sur $D = [0, +\infty[$.
\begin{itemize}
    \item \textbf{Limite simple :}
    Pour chaque $x \ge 0$ fixé, $\lim_{n \to +\infty} (n+x) = +\infty$, donc $\lim_{n \to +\infty} \frac{x}{n+x} = 0$.
    La limite simple est $f(x) = 0$.
    \item \textbf{Convergence uniforme :}
    Pour tout $n \ge 1$, la fonction $f_n$ est continue et croissante sur $[0, +\infty[$, et :
    \[ \lim_{x \to +\infty} f_n(x) = \lim_{x \to +\infty} \frac{x}{n+x} = 1 \]
    Par conséquent :
    \[ \|f_n - 0\|_{\infty, [0, +\infty[} = \sup_{x \ge 0} \frac{x}{n+x} = 1 \xrightarrow[n \to +\infty]{} 1 \neq 0 \]
    La suite $(f_n)$ \textbf{ne converge pas uniformément sur $[0, +\infty[$}.\\
    \textit{Suite test explicite :} Pour $x_n = n \in [0, +\infty[$, on a $f_n(x_n) = \frac{n}{n+n} = \frac{1}{2} \not\to 0$.
\end{itemize}

\textbf{Cas c) :} $f_n(x) = \frac{x^2}{n+x^2}$ sur $D = \mathbb{R}$.
\begin{itemize}
    \item \textbf{Limite simple :}
    Pour tout $x \in \mathbb{R}$ fixé, $\lim_{n \to +\infty} \frac{x^2}{n+x^2} = 0$.
    La limite simple est $f(x) = 0$.
    \item \textbf{Convergence uniforme :}
    Pour tout $n \ge 1$ :
    \[ \lim_{x \to +\infty} f_n(x) = \lim_{x \to +\infty} \frac{x^2}{n+x^2} = 1 \]
    D'où :
    \[ \|f_n - 0\|_{\infty, \mathbb{R}} = \sup_{x \in \mathbb{R}} \frac{x^2}{n+x^2} = 1 \xrightarrow[n \to +\infty]{} 1 \neq 0 \]
    La suite $(f_n)$ \textbf{ne converge pas uniformément sur $\mathbb{R}$}.\\
    \textit{Suite test explicite :} Pour $x_n = \sqrt{n} \in \mathbb{R}$, on a $f_n(x_n) = \frac{n}{n+n} = \frac{1}{2} \not\to 0$.
\end{itemize}
\end{proof}

\subsection*{Exercice 2: [Convergence uniforme et continuité ($f_n(x) = \frac{nx}{1+nx}$)]}

\textbf{Soru:}
On considère la suite de fonctions définies sur $[0, 1]$ par :
\[ f_n(x) = \frac{nx}{1+nx} \]
(1) Déterminer la limite simple $f$ de la suite $(f_n)$.\\
(2) La fonction $f$ est-elle continue sur $[0, 1]$ ?\\
(3) En déduire que la convergence n'est pas uniforme sur $[0, 1]$.\\
(4) Retrouver ce résultat en calculant $\sup_{x \in [0, 1]} |f_n(x) - f(x)|$.\\
(5) Montrer que la convergence est uniforme sur $[a, 1]$ pour tout $a \in ]0, 1[$.

\begin{proof}[Résolution]
\textbf{Partie (1) : Limite simple sur $[0, 1]$}
Soit $x \in [0, 1]$.
\begin{itemize}
    \item Si $x = 0$ : pour tout $n \ge 1$, $f_n(0) = \frac{0}{1+0} = 0 \to 0$.
    \item Si $x \in ]0, 1]$ : on a $nx \to +\infty$, d'où :
    \[ f_n(x) = \frac{nx}{1+nx} = \frac{1}{\frac{1}{nx} + 1} \xrightarrow[n \to +\infty]{} \frac{1}{0 + 1} = 1 \]
\end{itemize}
La suite $(f_n)$ converge simplement sur $[0, 1]$ vers la fonction :
\[ f(x) = \begin{cases} 0 & \text{si } x = 0 \\ 1 & \text{si } x \in ]0, 1] \end{cases} \]

\textbf{Partie (2) : Continuité de $f$}
Calculons la limite de $f$ en 0 à droite :
\[ \lim_{x \to 0^+} f(x) = 1 \]
Or $f(0) = 0$. Comme $\lim_{x \to 0^+} f(x) \neq f(0)$, la fonction $f$ est \textbf{discontinue au point $x = 0$}. Par conséquent, $f$ n'est pas continue sur $[0, 1]$.

\textbf{Partie (3) : Déduction de la non-uniformité}
Pour tout $n \ge 1$, la fonction $f_n$ est continue sur $[0, 1]$ (fonction rationnelle dont le dénominateur $1+nx \ge 1 > 0$ ne s'annule pas).
D'après le \textbf{théorème de transfert de la continuité}, si la suite de fonctions continues $(f_n)$ convergeait uniformément vers $f$ sur $[0, 1]$, alors la fonction limite $f$ serait continue sur $[0, 1]$.
Puisque $f$ n'est pas continue sur $[0, 1]$, par contraposée, la suite $(f_n)$ \textbf{ne converge pas uniformément vers $f$ sur $[0, 1]$}.

\textbf{Partie (4) : Calcul explicite du supremum}
Soit $x \in [0, 1]$. 
- En $x = 0$ : $|f_n(0) - f(0)| = |0 - 0| = 0$.\\
- Pour $x \in ]0, 1]$ :
\[ |f_n(x) - f(x)| = \left| \frac{nx}{1+nx} - 1 \right| = \left| \frac{nx - (1+nx)}{1+nx} \right| = \frac{1}{1+nx} \]
La fonction $x \mapsto \frac{1}{1+nx}$ est strictement décroissante sur $]0, 1]$. Par conséquent, sa borne supérieure est atteinte lorsque $x \to 0^+$ :
\[ \sup_{x \in [0, 1]} |f_n(x) - f(x)| = \sup_{x \in ]0, 1]} \frac{1}{1+nx} = \lim_{x \to 0^+} \frac{1}{1+nx} = 1 \]
Puisque $\|f_n - f\|_{\infty, [0, 1]} = 1 \xrightarrow[n \to +\infty]{} 1 \neq 0$, on retrouve bien que la convergence \textbf{n'est pas uniforme sur $[0, 1]$}.

\textbf{Partie (5) : Convergence uniforme sur $[a, 1]$ avec $a \in ]0, 1[$}
Sur le segment $[a, 1]$, on a $x \ge a > 0$, donc la limite simple est constante égale à 1 : $f(x) = 1$.
Pour tout $x \in [a, 1]$ :
\[ |f_n(x) - f(x)| = \frac{1}{1+nx} \]
Puisque la fonction $x \mapsto \frac{1}{1+nx}$ est strictement décroissante, son maximum sur $[a, 1]$ est atteint à la borne inférieure $x = a$ :
\[ \|f_n - f\|_{\infty, [a, 1]} = \sup_{x \in [a, 1]} \frac{1}{1+nx} = \frac{1}{1+na} \]
Puisque $a > 0$, on a :
\[ \lim_{n \to +\infty} \frac{1}{1+na} = 0 \]
D'où $\lim_{n \to +\infty} \|f_n - f\|_{\infty, [a, 1]} = 0$.
La convergence est donc \textbf{uniforme sur $[a, 1]$} pour tout $a \in ]0, 1[$.
\end{proof}

\subsection*{Exercice 3: [Suite de fonctions à maximum variable ($f_n(x) = nx(1-x)^n$)]}

\textbf{Soru:}
Pour tout $n \ge 1$, on pose :
\[ f_n(x) = nx(1-x)^n, \quad x \in [0, 1] \]
(1) Déterminer la limite simple de $(f_n)$ sur $[0, 1]$.\\
(2) Étudier les variations de $f_n$ et déterminer son maximum sur $[0, 1]$.\\
(3) La convergence est-elle uniforme sur $[0, 1]$ ?\\
(4) Montrer que la convergence est uniforme sur $[a, 1]$ pour tout $a \in ]0, 1[$.

\begin{proof}[Résolution]
\textbf{Partie (1) : Limite simple sur $[0, 1]$}
Soit $x \in [0, 1]$.
\begin{itemize}
    \item Si $x = 0$ : $f_n(0) = n \cdot 0 \cdot 1^n = 0 \to 0$.
    \item Si $x = 1$ : $f_n(1) = n \cdot 1 \cdot (1-1)^n = 0 \to 0$.
    \item Si $x \in ]0, 1[$ : on pose $q = 1-x$. Comme $x \in ]0, 1[$, on a $0 < q < 1$.
    D'après les croissances comparées usuelles, $\lim_{n \to +\infty} n q^n = 0$.
    Ainsi, $\lim_{n \to +\infty} f_n(x) = x \lim_{n \to +\infty} n(1-x)^n = 0$.
\end{itemize}
La suite $(f_n)$ converge simplement sur $[0, 1]$ vers la fonction nulle : $f(x) = 0$.

\textbf{Partie (2) : Variations et maximum de $f_n$}
Pour tout $n \ge 1$, la fonction $f_n$ est dérivable sur $[0, 1]$. Calculons sa dérivée :
\[ f_n'(x) = n(1-x)^n + nx \cdot n(1-x)^{n-1}(-1) = n(1-x)^{n-1} \left[ (1-x) - nx \right] = n(1-x)^{n-1} \left[ 1 - (n+1)x \right] \]
Pour $x \in ]0, 1[$, le terme $n(1-x)^{n-1}$ est strictement positif. Le signe de $f_n'(x)$ est donc exactement celui de $1 - (n+1)x$.
- $f_n'(x) > 0 \iff x < \frac{1}{n+1}$.\\
- $f_n'(x) = 0 \iff x = \frac{1}{n+1}$.\\
- $f_n'(x) < 0 \iff x > \frac{1}{n+1}$.\\
La fonction $f_n$ est donc strictement croissante sur $[0, \frac{1}{n+1}]$ et strictement décroissante sur $[\frac{1}{n+1}, 1]$.
Elle admet un maximum global absolu en $x_n = \frac{1}{n+1} \in [0, 1]$.
Calculons la valeur de ce maximum :
\[ M_n = f_n\left(\frac{1}{n+1}\right) = n \left(\frac{1}{n+1}\right) \left(1 - \frac{1}{n+1}\right)^n = \frac{n}{n+1} \left(\frac{n}{n+1}\right)^n = \left(\frac{n}{n+1}\right)^{n+1} = \left( 1 - \frac{1}{n+1} \right)^{n+1} \]
Calculons sa limite en utilisant la forme exponentielle :
\[ M_n = \exp\left( (n+1) \ln\left(1 - \frac{1}{n+1}\right) \right) \]
Or au voisinage de 0, $\ln(1 - u) = -u + o(u)$, donc :
\[ (n+1) \ln\left(1 - \frac{1}{n+1}\right) = (n+1) \left( -\frac{1}{n+1} + o\left(\frac{1}{n+1}\right) \right) = -1 + o(1) \]
Par continuité de l'exponentielle :
\[ \lim_{n \to +\infty} M_n = e^{-1} = \frac{1}{e} \]

\textbf{Partie (3) : Étude de la convergence uniforme sur $[0, 1]$}
On a :
\[ \|f_n - 0\|_{\infty, [0, 1]} = \sup_{x \in [0, 1]} |f_n(x)| = M_n = \left( 1 - \frac{1}{n+1} \right)^{n+1} \]
Par passage à la limite :
\[ \lim_{n \to +\infty} \|f_n - 0\|_{\infty, [0, 1]} = \frac{1}{e} \neq 0 \]
La suite $(f_n)$ \textbf{ne converge pas uniformément vers 0 sur $[0, 1]$}.
\textit{Remarque cruciale :} Les fonctions $f_n$ sont continues et la limite $f=0$ est continue. Pourtant, la convergence n'est pas uniforme ! C'est le phénomène de la bosse glissante qui s'échappe vers l'origine en conservant une hauteur finie ($1/e$).

\textbf{Partie (4) : Convergence uniforme sur $[a, 1]$ avec $a \in ]0, 1[$}
Soit $a \in ]0, 1[$. Puisque $\lim_{n \to +\infty} \frac{1}{n+1} = 0$, il existe un entier $N \in \mathbb{N}$ tel que pour tout $n \ge N$ :
\[ \frac{1}{n+1} < a \]
Ainsi, pour tout $n \ge N$, le point de maximum global $x_n = \frac{1}{n+1}$ se trouve strictement à gauche du segment $[a, 1]$.
Puisque $f_n$ est strictement décroissante sur $[\frac{1}{n+1}, 1]$, elle est en particulier strictement décroissante sur $[a, 1]$.
Par conséquent, son maximum sur $[a, 1]$ est atteint à la borne gauche $x = a$ :
\[ \sup_{x \in [a, 1]} |f_n(x)| = f_n(a) = na(1-a)^n \]
Comme $a \in ]0, 1[$, on a $0 < 1-a < 1$. Par croissances comparées :
\[ \lim_{n \to +\infty} na(1-a)^n = 0 \]
Il s'ensuit que $\lim_{n \to +\infty} \|f_n - 0\|_{\infty, [a, 1]} = 0$.
La convergence est donc \textbf{uniforme sur $[a, 1]$}.
\end{proof}

\subsection*{Exercice 4: [Critère de Cauchy uniforme (Suites et Séries de Fonctions)]}

\begin{tcolorbox}[colback=teal!5!white,colframe=teal!75!black,title=Hatırlatma: Düzgün Cauchy Kriteri (Critère de Cauchy Uniforme)]
Bir $(f_n)$ fonksiyon dizisinin $D$ üzerinde düzgün yakınsak olması için \textbf{gerek ve yeter şart}, dizinin $D$ üzerinde \textbf{düzgün Cauchy} olmasıdır:
\[ \forall \epsilon > 0, \exists N \in \mathbb{N}, \forall m > n \ge N, \forall x \in D : |f_m(x) - f_n(x)| < \epsilon \]
Bu kriter, limit fonksiyonu $f$'in ne olduğunu \textbf{bilmediğimiz veya hesaplayamadığımız} durumlarda düzgün yakınsaklığı kanıtlamak için kullanılan en güçlü araçtır!
\end{tcolorbox}

\textbf{Soru:}
(1) Soit $(f_n)$ une suite de fonctions définies sur un ensemble non-vide $D$. Montrer que la condition :
\[ \sup_{x \in D} |f_{n+1}(x) - f_n(x)| \xrightarrow[n \to +\infty]{} 0 \]
ne suffit pas à garantir que $(f_n)$ est uniformément de Cauchy sur $D$.\\
(2) Pour $n \ge 1$, on considère la suite de fonctions définies sur $[0, 1]$ par :
\[ f_n(x) = \sum_{k=1}^n \frac{x^k}{k^2} \]
Montrer, à l'aide du critère de Cauchy uniforme, que $(f_n)$ converge uniformément sur $[0, 1]$ sans chercher à déterminer sa limite.\\
(3) On considère la série de fonctions $\sum_{n=1}^{+\infty} \frac{x^n}{n}$ avec $x \in [0, 1[$. Montrer, à l'aide du critère de Cauchy uniforme, que cette série ne converge pas uniformément sur $[0, 1[$.

\begin{proof}[Résolution]
\textbf{Partie (1) : Insuffisance de la condition $|f_{n+1}-f_n| \to 0$}
Considérons la suite de fonctions constantes sur $D$ définie par les sommes partielles de la série harmonique :
\[ f_n(x) = \sum_{k=1}^n \frac{1}{k} = H_n \quad \text{pour tout } x \in D \]
Pour tout $x \in D$ :
\[ |f_{n+1}(x) - f_n(x)| = \frac{1}{n+1} \]
D'où $\sup_{x \in D} |f_{n+1}(x) - f_n(x)| = \frac{1}{n+1} \xrightarrow[n \to +\infty]{} 0$. La condition est donc bien satisfaite.
Pourtant, choisissons $m = 2n$. Alors pour tout $x \in D$ :
\[ |f_{2n}(x) - f_n(x)| = \sum_{k=n+1}^{2n} \frac{1}{k} \ge \sum_{k=n+1}^{2n} \frac{1}{2n} = n \cdot \frac{1}{2n} = \frac{1}{2} \]
En passant au supremum sur $D$, $\sup_{x \in D} |f_{2n}(x) - f_n(x)| \ge \frac{1}{2}$ pour tout $n \ge 1$.
Il est donc impossible de rendre cette quantité strictement inférieure à $\epsilon = \frac{1}{4}$ pour $n$ grand.
La suite $(f_n)$ n'est donc \textbf{pas uniformément de Cauchy sur $D$} (elle diverge d'ailleurs vers $+\infty$ en tout point).

\textbf{Partie (2) : Convergence uniforme de $f_n(x) = \sum_{k=1}^n \frac{x^k}{k^2}$ sur $[0, 1]$}
Soient $m, n \in \mathbb{N}$ tels que $m > n \ge 1$. Pour tout $x \in [0, 1]$ :
\[ |f_m(x) - f_n(x)| = \left| \sum_{k=n+1}^m \frac{x^k}{k^2} \right| = \sum_{k=n+1}^m \frac{x^k}{k^2} \]
Puisque $x \in [0, 1]$, on a $0 \le x^k \le 1$ pour tout $k \ge 1$. Par conséquent :
\[ |f_m(x) - f_n(x)| \le \sum_{k=n+1}^m \frac{1}{k^2} \le \sum_{k=n+1}^{+\infty} \frac{1}{k^2} \]
Cette majoration est totalement indépendante de $x \in [0, 1]$. En passant à la borne supérieure sur $[0, 1]$ :
\[ \sup_{x \in [0, 1]} |f_m(x) - f_n(x)| \le \sum_{k=n+1}^{+\infty} \frac{1}{k^2} \]
La série numérique $\sum_{k=1}^{+\infty} \frac{1}{k^2}$ est une série de Riemann convergente ($\alpha = 2 > 1$). Par propriété du reste d'une série convergente :
\[ R_n = \sum_{k=n+1}^{+\infty} \frac{1}{k^2} \xrightarrow[n \to +\infty]{} 0 \]
Soit $\epsilon > 0$. Il existe donc un entier $N \in \mathbb{N}$ tel que pour tout $n \ge N$, $R_n < \epsilon$.
Ainsi, pour tous $m > n \ge N$ :
\[ \sup_{x \in [0, 1]} |f_m(x) - f_n(x)| \le R_n < \epsilon \]
La suite $(f_n)$ est donc \textbf{uniformément de Cauchy sur $[0, 1]$}.
L'espace d'arrivée $\mathbb{R}$ muni de la valeur absolue étant complet, d'après le critère de Cauchy uniforme, la suite $(f_n)$ \textbf{converge uniformément sur $[0, 1]$}.

\textbf{Partie (3) : Non-convergence uniforme de $\sum_{n=1}^{+\infty} \frac{x^n}{n}$ sur $[0, 1[$}
Notons $S_n(x) = \sum_{k=1}^n \frac{x^k}{k}$ les sommes partielles de la série.
Supposons par l'absurde que la série converge uniformément sur $[0, 1[$.
D'après le critère de Cauchy uniforme, pour $\epsilon = \frac{1}{4} > 0$, il devrait exister un entier $N \in \mathbb{N}$ tel que pour tous $m > n \ge N$ :
\[ \sup_{x \in [0, 1[} |S_m(x) - S_n(x)| < \frac{1}{4} \]
Prenons en particulier $m = 2n$ pour un entier $n \ge N$. Pour tout $x \in [0, 1[$ :
\[ S_{2n}(x) - S_n(x) = \sum_{k=n+1}^{2n} \frac{x^k}{k} \]
Puisque chaque terme $x \mapsto \frac{x^k}{k}$ est continu et croissant sur $[0, 1[$, on a :
\[ \sup_{x \in [0, 1[} |S_{2n}(x) - S_n(x)| = \lim_{x \to 1^-} \sum_{k=n+1}^{2n} \frac{x^k}{k} = \sum_{k=n+1}^{2n} \frac{1}{k} \]
Or, par minoration usuelle :
\[ \sum_{k=n+1}^{2n} \frac{1}{k} \ge n \cdot \frac{1}{2n} = \frac{1}{2} \]
On obtient donc :
\[ \sup_{x \in [0, 1[} |S_{2n}(x) - S_n(x)| \ge \frac{1}{2} > \frac{1}{4} \]
ce qui contredit l'hypothèse de Cauchy uniforme.
Par conséquent, la série de fonctions $\sum_{n=1}^{+\infty} \frac{x^n}{n}$ \textbf{ne converge pas uniformément sur $[0, 1[$}.
\end{proof}

\subsection*{Exercice 5: [Séries de fonctions (Convergence Simple, Absolue, Uniforme et Normale)]}

\textbf{Soru:}
Étudier la convergence simple, absolue, uniforme et normale des séries suivantes sur les domaines indiqués :
\[ a) \sum_{n=1}^{+\infty} \frac{\sin(nx)}{n^2}, x \in \mathbb{R} \]
\[ b) \sum_{n=1}^{+\infty} \frac{(-1)^n x}{n}, x \in [0, 1] \]
\[ c) \sum_{n=1}^{+\infty} \frac{x}{(1+nx)(1+(n+1)x)}, x \in [0, +\infty[ \]
\textbf{Question théorique :} Soit $\sum_{n=1}^{+\infty} f_n$ une série de fonctions convergeant simplement sur un ensemble $D$, de somme $S$. On note $R_N(x) = S(x) - \sum_{n=1}^N f_n(x)$.
Montrer que la série converge uniformément sur $D$ si et seulement si $\sup_{x \in D} |R_N(x)| \xrightarrow[N \to +\infty]{} 0$.

\begin{proof}[Résolution]
\textbf{Série a) :} $\sum_{n=1}^{+\infty} u_n(x)$ avec $u_n(x) = \frac{\sin(nx)}{n^2}$ sur $D = \mathbb{R}$.
\begin{itemize}
    \item \textbf{Convergence normale :}
    Pour tout $x \in \mathbb{R}$, on sait que $|\sin(nx)| \le 1$. Ainsi :
    \[ |u_n(x)| = \left| \frac{\sin(nx)}{n^2} \right| \le \frac{1}{n^2} \]
    Cette borne est indépendante de $x$. En passant à la borne supérieure :
    \[ \|u_n\|_{\infty, \mathbb{R}} = \sup_{x \in \mathbb{R}} \left| \frac{\sin(nx)}{n^2} \right| \le \frac{1}{n^2} \]
    (et cette valeur est effectivement atteinte pour $x = \frac{\pi}{2n}$, donc $\|u_n\|_{\infty, \mathbb{R}} = \frac{1}{n^2}$).\\
    La série numérique $\sum \frac{1}{n^2}$ est une série de Riemann convergente ($\alpha = 2 > 1$).
    Par le \textbf{critère de Weierstrass}, la série $\sum u_n$ \textbf{converge normalement sur $\mathbb{R}$}.
    \item \textbf{Autres convergences :}
    Puisque la convergence normale implique toutes les autres formes de convergence, la série converge également \textbf{uniformément}, \textbf{absolument} et \textbf{simplement sur $\mathbb{R}$}.
\end{itemize}

\textbf{Série b) :} $\sum_{n=1}^{+\infty} u_n(x)$ avec $u_n(x) = \frac{(-1)^n x}{n}$ sur $D = [0, 1]$.
\begin{itemize}
    \item \textbf{Convergence simple :}
    - Pour $x = 0$ : $u_n(0) = 0$ pour tout $n$, la somme vaut 0.\\
    - Pour $x \in ]0, 1]$ : on a $\sum \frac{(-1)^n x}{n} = x \sum \frac{(-1)^n}{n}$. C'est le produit de $x$ par la série harmonique alternée, qui converge d'après le critère de Leibniz.\\
    Donc la série \textbf{converge simplement sur $[0, 1]$}.
    \item \textbf{Convergence absolue :}
    Pour tout $x \in ]0, 1]$, $|u_n(x)| = \frac{x}{n}$. La série $\sum \frac{x}{n} = x \sum \frac{1}{n}$ diverge car $\sum \frac{1}{n}$ est la série harmonique divergente.
    La série \textbf{ne converge absolument en aucun point de $]0, 1]$} (elle ne converge absolument qu'au point isolé $x = 0$).
    \item \textbf{Convergence normale :}
    $\|u_n\|_{\infty, [0, 1]} = \sup_{x \in [0, 1]} \frac{x}{n} = \frac{1}{n}$.
    La série numérique associée est $\sum \frac{1}{n}$, qui diverge.
    Par conséquent, la série \textbf{ne converge pas normalement sur $[0, 1]$}.
    \item \textbf{Convergence uniforme :}
    Pour chaque $x \in [0, 1]$, la série est alternée et satisfait les hypothèses du critère de Leibniz. D'après la formule de majoration du reste des séries alternées :
    \[ |R_N(x)| = \left| \sum_{n=N+1}^{+\infty} \frac{(-1)^n x}{n} \right| \le |u_{N+1}(x)| = \frac{x}{N+1} \]
    Comme $x \in [0, 1]$, on a immédiatement :
    \[ |R_N(x)| \le \frac{1}{N+1} \]
    Cette majoration est indépendante de $x$. En passant au supremum sur $[0, 1]$ :
    \[ \sup_{x \in [0, 1]} |R_N(x)| \le \frac{1}{N+1} \xrightarrow[N \to +\infty]{} 0 \]
    La série \textbf{converge uniformément sur $[0, 1]$}.\\
    \textit{Remarque théorique majeure :} Cet exercice démontre qu'une série peut être uniformément convergente tout en n'étant ni normalement convergente, ni absolument convergente !
\end{itemize}

\textbf{Série c) :} $\sum_{n=1}^{+\infty} u_n(x)$ avec $u_n(x) = \frac{x}{(1+nx)(1+(n+1)x)}$ sur $D = [0, +\infty[$.
\begin{itemize}
    \item \textbf{Décomposition télescopique :}
    Calculons la différence :
    \[ \frac{1}{1+nx} - \frac{1}{1+(n+1)x} = \frac{(1+(n+1)x) - (1+nx)}{(1+nx)(1+(n+1)x)} = \frac{x}{(1+nx)(1+(n+1)x)} = u_n(x) \]
    Il s'agit donc d'une \textbf{série télescopique} !
    \item \textbf{Sommes partielles et convergence simple :}
    Pour tout $x \ge 0$, la somme partielle d'ordre $N$ s'écrit par télescopage :
    \[ S_N(x) = \sum_{n=1}^N u_n(x) = \sum_{n=1}^N \left( \frac{1}{1+nx} - \frac{1}{1+(n+1)x} \right) = \frac{1}{1+x} - \frac{1}{1+(N+1)x} \]
    Étudions la limite quand $N \to +\infty$ selon la valeur de $x$ :
    - Si $x = 0$ : $u_n(0) = 0$ pour tout $n$, donc $S_N(0) = 0 \to 0$.\\
    - Si $x > 0$ : $\lim_{N \to +\infty} \frac{1}{1+(N+1)x} = 0$, d'où $\lim_{N \to +\infty} S_N(x) = \frac{1}{1+x}$.\\
    La série converge simplement sur $[0, +\infty[$ vers la fonction somme :
    \[ S(x) = \begin{cases} 0 & \text{si } x = 0 \\ \frac{1}{1+x} & \text{si } x > 0 \end{cases} \]
    \item \textbf{Convergence uniforme :}
    Pour tout $n \ge 1$, la fonction $u_n$ est continue sur $[0, +\infty[$, donc chaque somme partielle $S_N$ est continue sur $[0, +\infty[$.
    Observons la limite de la fonction somme $S$ en 0 :
    \[ \lim_{x \to 0^+} S(x) = \lim_{x \to 0^+} \frac{1}{1+x} = 1 \neq S(0) = 0 \]
    La fonction somme $S$ est discontinue en 0.
    Par le \textbf{théorème de continuité des séries de fonctions}, si la série convergeait uniformément sur $[0, +\infty[$, sa somme $S$ serait continue sur cet intervalle.
    Par contraposée, la série \textbf{ne converge pas uniformément sur $[0, +\infty[$}.
    \item \textbf{Convergence absolue et normale :}
    Comme $u_n(x) \ge 0$ pour tout $x \ge 0$, la convergence absolue est équivalente à la convergence simple : elle converge donc absolument sur $[0, +\infty[$.
    En revanche, elle ne converge pas normalement sur $[0, +\infty[$ (car la convergence normale impliquerait la convergence uniforme, ce qui est faux).
\end{itemize}

\textbf{Démonstration de la question théorique :}
Par définition, la série $\sum_{n=1}^{+\infty} f_n$ converge uniformément sur $D$ vers sa somme $S$ si et seulement si la suite des sommes partielles $S_N = \sum_{n=1}^N f_n$ converge uniformément vers $S$ sur $D$.
Par définition de la convergence uniforme d'une suite de fonctions, cela équivaut à :
\[ \lim_{N \to +\infty} \|S - S_N\|_{\infty, D} = 0 \iff \lim_{N \to +\infty} \sup_{x \in D} |S(x) - S_N(x)| = 0 \]
Or, par définition du reste d'ordre $N$ :
\[ R_N(x) = S(x) - \sum_{n=1}^N f_n(x) = S(x) - S_N(x) \]
On a donc identiquement $|S(x) - S_N(x)| = |R_N(x)|$.
Par conséquent, la série converge uniformément sur $D$ si et seulement si :
\[ \lim_{N \to +\infty} \sup_{x \in D} |R_N(x)| = 0 \]
Ce qui achève la démonstration.
\end{proof}

\subsection*{Exercice 6: [Vrai ou Faux: 8 Propositions Fondamentales sur les Suites et Séries de Fonctions]}

\textbf{Soru:}
Pour chacune des affirmations suivantes, déterminer si elle est vraie ou fausse. Justifier toute réponse affirmative par une démonstration et toute réponse négative par un contre-exemple.
\begin{enumerate}
    \item Si une suite de fonctions continues converge simplement alors sa limite est continue.
    \item Si une suite de fonctions converge simplement vers une fonction continue alors la convergence est uniforme.
    \item Si une suite de fonctions converge uniformément sur tout segment de $]0, 1[$ alors elle converge uniformément sur $]0, 1[$.
    \item Si une suite de fonctions converge uniformément sur $D$ alors elle converge uniformément sur toute partie $A \subset D$.
    \item Si une série de fonctions converge absolument en tout point de $D$ alors elle converge uniformément sur $D$.
    \item Si une série de fonctions converge normalement sur $D$ alors elle converge uniformément sur $D$.
    \item Si une série de fonctions converge uniformément sur $D$ alors elle converge normalement sur $D$.
    \item Si une série $\sum f_n$ converge uniformément sur $D$ alors la suite $(f_n)$ converge uniformément vers zéro sur $D$.
\end{enumerate}

\begin{proof}[Résolution]
\textbf{Proposition 1 : "Si une suite de fonctions continues converge simplement alors sa limite est continue."}
\begin{center}
\textbf{\textcolor{red}{FAUX}}
\end{center}
\textbf{Contre-exemple :} Posons $f_n(x) = x^n$ sur $D = [0, 1]$.
Chaque fonction $f_n$ est continue sur $[0, 1]$. Pourtant, sa limite simple est :
\[ f(x) = \begin{cases} 0 & \text{si } x \in [0, 1[ \\ 1 & \text{si } x = 1 \end{cases} \]
qui est discontinue en $x = 1$. La convergence simple ne préserve pas la continuité.

\vspace{0.3cm}
\textbf{Proposition 2 : "Si une suite de fonctions converge simplement vers une fonction continue alors la convergence est uniforme."}
\begin{center}
\textbf{\textcolor{red}{FAUX}}
\end{center}
\textbf{Contre-exemple :} Posons $f_n(x) = \frac{nx}{1 + n^2 x^2}$ sur $D = [0, 1]$.
Pour tout $x \in [0, 1]$, la limite simple est la fonction nulle $f(x) = 0$, qui est parfaitement continue sur $[0, 1]$.
Pourtant, en $x_n = \frac{1}{n} \in [0, 1]$, on a $f_n(x_n) = \frac{1}{2}$, d'où $\|f_n - 0\|_{\infty, [0, 1]} = \frac{1}{2} \not\to 0$.
La convergence n'est pas uniforme. La continuité de la limite est une condition nécessaire mais nullement suffisante.

\vspace{0.3cm}
\textbf{Proposition 3 : "Si une suite de fonctions converge uniformément sur tout segment de $]0, 1[$ alors elle converge uniformément sur $]0, 1[$."}
\begin{center}
\textbf{\textcolor{red}{FAUX}}
\end{center}
\textbf{Contre-exemple :} Posons $f_n(x) = x^n$ sur $D = ]0, 1[$.
- Soit $[a, b] \subset ]0, 1[$ un segment arbitraire, avec $0 < a \le b < 1$. Puisque $x \mapsto x^n$ est croissante :
\[ \sup_{x \in [a, b]} |f_n(x) - 0| = b^n \xrightarrow[n \to +\infty]{} 0 \quad (\text{car } b < 1) \]
La suite converge donc uniformément sur tout segment de $]0, 1[$.
- En revanche, sur l'intervalle ouvert $]0, 1[$ tout entier :
\[ \sup_{x \in ]0, 1[} |f_n(x) - 0| = \sup_{x \in ]0, 1[} x^n = 1 \xrightarrow[n \to +\infty]{} 1 \neq 0 \]
La convergence n'est pas uniforme sur $]0, 1[$.

\vspace{0.3cm}
\textbf{Proposition 4 : "Si une suite de fonctions converge uniformément sur $D$ alors elle converge uniformément sur toute partie $A \subset D$."}
\begin{center}
\textbf{\textcolor{blue!80!black}{VRAI}}
\end{center}
\textbf{Démonstration :} Soit $A \subset D$. Par propriété de la borne supérieure sur une partie :
\[ \sup_{x \in A} |f_n(x) - f(x)| \le \sup_{x \in D} |f_n(x) - f(x)| \]
Puisque $(f_n)$ converge uniformément vers $f$ sur $D$, par hypothèse :
\[ \lim_{n \to +\infty} \sup_{x \in D} |f_n(x) - f(x)| = 0 \]
Comme $0 \le \sup_{x \in A} |f_n(x) - f(x)|$, par le théorème d'encadrement des limites, on obtient :
\[ \lim_{n \to +\infty} \sup_{x \in A} |f_n(x) - f(x)| = 0 \]
Ce qui établit la convergence uniforme sur $A$.

\vspace{0.3cm}
\textbf{Proposition 5 : "Si une série de fonctions converge absolument en tout point de $D$ alors elle converge uniformément sur $D$."}
\begin{center}
\textbf{\textcolor{red}{FAUX}}
\end{center}
\textbf{Contre-exemple :} Considérons la série géométrique $\sum_{n=0}^{+\infty} x^n$ sur $D = [0, 1[$.
Pour chaque $x \in [0, 1[$, on a $|x^n| = x^n$ et la série numérique $\sum x^n = \frac{1}{1-x}$ converge. La série converge donc absolument en chaque point de $D$.
Pourtant, le reste d'ordre $N$ s'écrit $R_N(x) = \sum_{n=N+1}^{+\infty} x^n = \frac{x^{N+1}}{1-x}$.
Pour tout $N \in \mathbb{N}$ fixé :
\[ \lim_{x \to 1^-} R_N(x) = \lim_{x \to 1^-} \frac{x^{N+1}}{1-x} = +\infty \]
Donc $\sup_{x \in [0, 1[} |R_N(x)| = +\infty$. La série ne converge donc absolument pas uniformément sur $[0, 1[$.

\vspace{0.3cm}
\textbf{Proposition 6 : "Si une série de fonctions converge normalement sur $D$ alors elle converge uniformément sur $D$."}
\begin{center}
\textbf{\textcolor{blue!80!black}{VRAI}}
\end{center}
\textbf{Démonstration :} Par hypothèse de convergence normale, la série numérique à termes positifs $\sum_{n=1}^{+\infty} \|f_n\|_{\infty, D}$ converge.
Soit $\epsilon > 0$. D'après le critère de Cauchy pour les séries numériques, il existe $N \in \mathbb{N}$ tel que pour tous $m > n \ge N$ :
\[ \sum_{k=n+1}^m \|f_k\|_{\infty, D} < \epsilon \]
Or, pour tout $x \in D$, par l'inégalité triangulaire :
\[ \left| \sum_{k=n+1}^m f_k(x) \right| \le \sum_{k=n+1}^m |f_k(x)| \le \sum_{k=n+1}^m \|f_k\|_{\infty, D} \]
En passant à la borne supérieure sur $D$ :
\[ \sup_{x \in D} \left| \sum_{k=n+1}^m f_k(x) \right| \le \sum_{k=n+1}^m \|f_k\|_{\infty, D} < \epsilon \]
La suite des sommes partielles vérifie donc le \textbf{critère de Cauchy uniforme} sur $D$.
L'espace $\mathbb{R}$ étant complet, la série $\sum f_n$ converge uniformément sur $D$.

\vspace{0.3cm}
\textbf{Proposition 7 : "Si une série de fonctions converge uniformément sur $D$ alors elle converge normalement sur $D$."}
\begin{center}
\textbf{\textcolor{red}{FAUX}}
\end{center}
\textbf{Contre-exemple :} Considérons la série $\sum_{n=1}^{+\infty} \frac{(-1)^n x}{n}$ sur $D = [0, 1]$ (étudiée à l'Exercice 5b).
- Par le critère de Leibniz, $|R_N(x)| \le \frac{x}{N+1} \le \frac{1}{N+1}$, d'où $\sup_{x \in [0, 1]} |R_N(x)| \le \frac{1}{N+1} \to 0$. La série converge donc uniformément sur $[0, 1]$.
- Pourtant, $\|f_n\|_{\infty, [0, 1]} = \sup_{x \in [0, 1]} \frac{x}{n} = \frac{1}{n}$. La série numérique $\sum \frac{1}{n}$ est la série harmonique divergente. La série ne converge donc pas normalement.
(Autre exemple encore plus simple : la série constante $f_n(x) = \frac{(-1)^n}{n}$ sur $[0, 1]$).

\vspace{0.3cm}
\textbf{Proposition 8 : "Si une série $\sum f_n$ converge uniformément sur $D$ alors la suite $(f_n)$ converge uniformément vers zéro sur $D$."}
\begin{center}
\textbf{\textcolor{blue!80!black}{VRAI}}
\end{center}
\textbf{Démonstration :} Notons $S_n = \sum_{k=1}^n f_k$ la suite des sommes partielles et $S$ la somme de la série.
Par hypothèse, la série converge uniformément vers $S$ sur $D$, ce qui signifie que :
\[ \|S_n - S\|_{\infty, D} \xrightarrow[n \to +\infty]{} 0 \]
Pour tout $n \ge 2$, on a la relation :
\[ f_n = S_n - S_{n-1} = (S_n - S) - (S_{n-1} - S) \]
Par l'inégalité triangulaire pour la norme uniforme (ou pour la valeur absolue en tout point $x \in D$) :
\[ \|f_n\|_{\infty, D} = \|(S_n - S) - (S_{n-1} - S)\|_{\infty, D} \le \|S_n - S\|_{\infty, D} + \|S_{n-1} - S\|_{\infty, D} \]
Puisque $\lim_{n \to +\infty} \|S_n - S\|_{\infty, D} = 0$ et $\lim_{n \to +\infty} \|S_{n-1} - S\|_{\infty, D} = 0$, la somme des deux termes tend vers 0.
Par encadrement :
\[ \lim_{n \to +\infty} \|f_n - 0\|_{\infty, D} = 0 \]
Ce qui prouve que la suite $(f_n)$ \textbf{converge uniformément vers la fonction nulle sur $D$}.
\end{proof}

\end{document}
